\documentclass[11pt,a4paper]{article}
\usepackage[utf8]{inputenc}
\usepackage[T1]{fontenc}
\usepackage{amsmath,amssymb,amsthm}
\usepackage{geometry}
\usepackage{fancyhdr}
\usepackage{titlesec}
\usepackage{enumitem}
\usepackage{array}
\usepackage{booktabs}

\geometry{margin=1in}

\pagestyle{fancy}
\fancyhf{}
\fancyhead[C]{\small Cayley, Collected Mathematical Papers, Vol.\ I}
\fancyfoot[C]{\thepage}
\renewcommand{\headrulewidth}{0.4pt}

\theoremstyle{definition}
\newtheorem*{lemma}{Lemma}
\newtheorem*{problem}{Problem}

\titleformat{\section}{\large\bfseries}{\thesection.}{0.5em}{}
\titleformat{\subsection}{\normalsize\bfseries}{\thesubsection.}{0.5em}{}

\newcommand{\mb}[1]{\mathbf{#1}}
\DeclareMathOperator{\lgn}{\lambda}

\begin{document}

\begin{center}
\Large\textbf{Collected Mathematical Papers of Arthur Cayley, Vol.\ I}\\[6pt]
\large\textit{Pages 417--468}
\end{center}

\vspace{1em}

% ============================================================
% PAPER 70 (conclusion)
% ============================================================
\section*{[70. Sur quelques th\'{e}or\`{e}mes de la g\'{e}om\'{e}trie de position.\ (concluded)]}
\markboth{Sur quelques th\'{e}or\`{e}mes de la g\'{e}om\'{e}trie de position.}{}
\addcontentsline{toc}{section}{70. Sur quelques th\'{e}or\`{e}mes de la g\'{e}om\'{e}trie de position (concluded)}

\begin{flushright}
\textit{[From the Journal f"{u}r die reine und angewandte Mathematik (Crelle),\ tom.\ xxxviii.\ (1848),\\
pp.\ 413--419.]}
\end{flushright}

\setcounter{page}{417}

plans sera \'{e}gal au rapport anharmonique des quatre points; ce qui suffit pour d\'{e}terminer le point $Q$, puisque les quatre plans et les trois autres points sont donn\'{e}s, et la construction graphique pour d\'{e}terminer ce point $Q$ est parfaitement connue. Il est d'ailleurs \'{e}vident que la droite men\'{e}e par un point donn\'{e}, de sorte qu'elle rencontre des droites r\'{e}ciproques gauches l'une \`{a} l'autre, est situ\'{e}e dans la r\'{e}ciproque gauche de ce point. Donc, en menant par le point $q$ la droite qui rencontre les deux droites $AC$ et $BD$, le plan passant par cette droite et par le point $Q$, est la r\'{e}ciproque gauche du point $q$ qu'il s'agissait de trouver\textsuperscript{1}. \'{E}galement, on pourrait construire la r\'{e}ciproque gauche d'un plan donn\'{e}.

Donc enfin, pour trouver les r\'{e}ciproques d'un point de la premi\`{e}re surface:

\begin{enumerate}[label=(\Alph*)]
\item ``Construisez le plan tangent \`{a} ce point de la premi\`{e}re surface, et construisez de la mani\`{e}re expliqu\'{e}e ci-dessus, la r\'{e}ciproque gauche du point. Par la droite d'intersection de ces deux plans menez deux plans tangents \`{a} la seconde surface: ces deux plans seront les r\'{e}ciproques qu'il s'agissait de trouver.''

\item En construisant les r\'{e}ciproques du point de contact avec la premi\`{e}re surface du plan dont il s'agit, les points de contact avec la seconde surface, de ces deux plans, seront les r\'{e}ciproques dont il s'agissait.

\item ``Construisez le point de contact de ce plan avec la seconde surface, et construisez la r\'{e}ciproque gauche de ce m\^{e}me plan: la droite men\'{e}e par ces deux points rencontrera la premi\`{e}re surface dans deux points qui seront les r\'{e}ciproques que l'on d\'{e}sirait.''

\item ``Construisez les r\'{e}ciproques du plan tangent passant par le point donn\'{e} de la seconde surface: les plans tangents \`{a} la premi\`{e}re surface passant par ces deux points, seront les r\'{e}ciproques qu'il s'agissait de trouver.''
\end{enumerate}

En effet les th\'{e}or\`{e}mes (C, D) ne sont que des transformations des th\'{e}or\`{e}mes (A, B) au moyen de la th\'{e}orie des polaires r\'{e}ciproques.

Enfin, pour trouver les r\'{e}ciproques d'un point quelconque, menez par ce point trois plans tangents ou \`{a} la premi\`{e}re ou \`{a} la seconde surface, et construisez les r\'{e}ciproques de ces plans au moyen du th\'{e}or\`{e}me (B ou C): les deux plans men\'{e}s par ces points r\'{e}ciproques, pris trois et trois ensemble, et combin\'{e}s de mani\`{e}re que l'intersection des deux plans coincide avec l'intersection de la polaire par rapport \`{a} la premi\`{e}re surface et de la r\'{e}ciproque gauche du point donn\'{e}, selon la remarque ci-dessus, seront les r\'{e}ciproques cherch\'{e}es; et de la m\^{e}me mani\`{e}re pour les r\'{e}ciproques d'un plan quelconque.

\vspace{0.5em}
\noindent\rule{0.3\textwidth}{0.4pt}

\small
\textsuperscript{1} Il y a un cas tr\`{e}s simple qui m\'{e}rite d'\^{e}tre consid\'{e}r\'{e}; savoir celui o\`{u} le point $P$ co\"{i}ncide avec $A$. Dans ce cas $Q$ co\"{i}ncide aussi avec $A$, et la r\'{e}ciproque gauche de $q$ se r\'{e}duit au plan $qAC$. De m\^{e}me, si les points $P$, $B$ co\"{i}ncident, la r\'{e}ciproque gauche de $q$ se r\'{e}duit au plan $qBD$.
\normalsize

\newpage

\setcounter{page}{418}

Je n'essaierai pas d'\'{e}num\'{e}rer ici le grand nombre de relations descriptives qui pourraient \^{e}tre tir\'{e}es des constructions qu'on vient d'expliquer. L'on remarquera sans peine l'analogie parfaite qui existe pour toute cette th\'{e}orie et la th\'{e}orie correspondante de la g\'{e}om\'{e}trie \`{a} deux dimensions, telle que M.~Pl"{u}cker l'a expos\'{e}e, ``System der analytischen Geometrie,'' [4$^{\circ}$, Berlin, 1835], pp.~78--83. On pourra aussi consulter sur ce point un m\'{e}moire [61] que je viens de composer pour le Cambridge and Dublin Mathematical Journal, et qui para\^{i}tra prochainement.

La v\'{e}rification analytique de ces th\'{e}or\`{e}mes donne lieu \`{a} des d\'{e}veloppements assez int\'{e}ressants. Pour obtenir l'\'{e}quation de la premi\`{e}re des surfaces du second ordre, il suffit d'\'{e}crire $\xi$, $\eta$, $\zeta$, $\omega$ au lieu de $x$, $y$, $z$, $w$, dans l'\'{e}quation de l'un ou de l'autre des plans r\'{e}ciproques. En supposant que $x=0$, $y=0$, $z=0$, $w=0$ soient des plans conjugu\'{e}s par rapport \`{a} la surface, et en remarquant que chacune de ces coordonn\'{e}es peut \^{e}tre cens\'{e}e contenir un facteur constant, l'\'{e}quation de la premi\`{e}re surface peut \^{e}tre \'{e}crite sous la forme
\[
x^{2} + y^{2} + z^{2} + w^{2} = 0.
\]
On aura alors pour $a$, $b$, $c$, $d$; $a'$, $b'$, $c'$, $d'$; $a''$, $b''$, $c''$, $d''$; $a'''$, $b'''$, $c'''$, $d'''$ un syst\`{e}me de la forme $l$, $-h$, $g$, $-a$; $h$, $l$, $-f$, $-b$; $-g$, $f$, $l$, $-c$; $a$, $b$, $c$, $l$, et les \'{e}quations des plans r\'{e}ciproques deviendront
\begin{align*}
&(\xi x + \eta y + \zeta z + \omega w) \pm x(-h\eta + g\zeta - a\omega) \pm y(h\xi - f\zeta - b\omega) \\
&\qquad \pm z(-g\xi + f\eta - c\omega) \pm w(a\xi + b\eta + c\zeta) = 0,
\end{align*}
ce qui prouve le th\'{e}or\`{e}me \'{e}nonc\'{e} ci-dessus; savoir que les deux r\'{e}ciproques passent par la droite d'intersection de la polaire et de la r\'{e}ciproque gauche du point. On obtient sans difficult\'{e}, pour \'{e}quation de la seconde surface:
\[
(x - hy + gz - aw)^{2} + (hx + y - fz - bw)^{2} + (-gx + fy + z - cw)^{2} + (ax + by + cz + w)^{2} = 0,
\]
ou sous une autre forme plus commode:
\begin{multline*}
(x^{2} + y^{2} + z^{2} + w^{2}) \\
+ (-hy + gz - aw)^{2} + (hx - fz - bw)^{2} + (-gx + fy - cw)^{2} + (ax + by + cz)^{2} = 0.
\end{multline*}
Les coordonn\'{e}es des points $A$, $B$, $C$, $D$ seront d\'{e}termin\'{e}es au moyen des expressions
\[
\frac{-hy + gz - aw}{x} = \frac{hx - fz - bw}{y} = \frac{-gx + fy - cw}{z} = \frac{ax + by + cz}{w},
\]
ou, en introduisant la quantit\'{e} ind\'{e}termin\'{e}e $s$:
\begin{align*}
sx - hy + gz - aw &= 0, \\
hx + sy - fz - bw &= 0, \\
-gx + fy + sz - cw &= 0, \\
ax + by + cz + sw &= 0.
\end{align*}

\newpage

\setcounter{page}{419}

En effet, on d\'{e}montrera ais\'{e}ment, non seulement que les points d\'{e}finis par ces \'{E}quations sont situ\'{e}s dans l'intersection des deux surfaces, mais aussi que les deux surfaces se touchent dans ces points-ci; ce qui fait voir qu'en combinant convenablement les quatre points, les deux surfaces doivent se couper (comme nous l'avons d\'{e}j\`{a} avanc\'{e}), suivant les droites $AB$, $BC$, $CD$, $DA$.

Je vais examiner encore de plus pr\`{e}s le syst\`{e}me d'\'{e}quations qui d\'{e}terminent ces points. On en tire tout de suite, pour d\'{e}terminer $s$, l'\'{e}quation
\[
s^{4} + \Pi s^{2} + \Theta^{2} = 0,
\]
o\`{u} l'on a fait pour abr\'{e}ger:
\[
\Pi = a^{2} + b^{2} + c^{2} + l^{2} + f^{2} + g^{2} + h^{2}, \qquad \Theta = af + bg + ch.
\]
Supposons encore
\[
s^{2}a + f\Theta = F, \quad s^{2}b + g\Theta = G, \quad s^{2}c + h\Theta = H;
\]
l'\'{e}quation qui d\'{e}termine $s$, pourra \^{e}tre \'{e}crite sous cette autre forme:
\[
AF + BG + CH = 0.
\]

J'ai trouv\'{e} que les \'{E}quations des droites $AC$ et $BD$ peuvent \^{e}tre pr\'{e}sent\'{e}es sous les formes assez \'{e}l\'{e}gantes
\begin{align*}
-Hy + Gz - Aw &= 0, \qquad -Cy + Bz - Fw = 0, \\
Hx \qquad - Fz - Bw &= 0, \qquad Cx \qquad - Az - Gw = 0, \\
-Gx + Fy \qquad - Cw &= 0, \qquad -Bx + Ay \qquad - Hw = 0, \\
Ax + By + Cz \qquad\qquad &= 0, \qquad Fx + Gy + Hz \qquad\qquad = 0,
\end{align*}
ou chacun de ces syst\`{e}mes d'\'{e}quations n'est \'{e}quivalent qu'\`{a} un syst\`{e}me de deux \'{E}quations. En entrechangeant les racines de l'\'{e}quation en $s^{2}$, c'est-\`{a}-dire en \'{e}crivant $-s^{2}$ au lieu de $s^{2}$, les deux syst\`{e}mes ne font que s'entrechanger; comme en effet cela doit \^{e}tre.

Enfin, les \'{e}quations des plans $ACD$, $BAC$, ou des plans $ABD$, $BDC$ peuvent \^{e}tre exprim\'{e}es sous les formes
\begin{align*}
PQ &= (-Hy + Gz - Aw)^{2} + (Hx - Fz - Bw)^{2} + (-Gx + Fy - Cw)^{2} + (Ax + By + Cz)^{2} = 0, \\
RS &= (-Cy + Bz - Fw)^{2} + (Cx - Az - Gw)^{2} + (-Bx + Ay - Hw)^{2} + (Fx + Gy + Hz)^{2} = 0;
\end{align*}
\'{E}quations desquelles on tire les relations identiques
\begin{align*}
PQ + RS &= -8s^{2}(\Pi^{2} - 4\Theta^{2})(x^{2} + y^{2} + z^{2} + w^{2}), \\
s^{2}PQ + 8^{2}RS &= 8^{2}(\Pi^{2} - 4\Theta^{2}) \\
&\quad \times [(-hy + gz - aw)^{2} + (hx - fz - bw)^{2} + (-gx + fy - cw)^{2} + (ax + by + cz)^{2}],
\end{align*}
qui mettent en \'{e}vidence ce que l'on savait d\'{e}j\`{a}, savoir, que les deux surfaces contiennent les droites $AB$, $BC$, $CD$, $DA$.

\newpage

\setcounter{page}{420}

Mettons pour un moment, pour abr\'{e}ger:
\begin{align*}
\phantom{-}l\xi - h\eta + g\zeta - a\omega &= L, \\
-h\xi + l\eta - f\zeta - b\omega &= M, \\
-g\xi + f\eta + l\zeta - c\omega &= N, \\
a\xi + b\eta + c\zeta + l\omega &= P,
\end{align*}
de sorte que l'\'{e}quation de la r\'{e}ciproque gauche du point $(\xi, \eta, \zeta, \omega)$ devient
\[
Lx + My + Nz + Pw = 0;
\]
supposons de plus que, en combinant cette \'{e}quation avec les \'{e}quations des droites $AC$ et $BD$ respectivement, on obtient $x:y:z:w = X:Y:Z:W$ et $x:y:z:w = X':Y':Z':W'$ respectivement. De l\`{a} on tire
\begin{align*}
X &= \phantom{-}CM + BN - FP, &\quad X' &= \phantom{-}HM + GN - AP, \\
Y &= -CL \qquad - AN - GP, &\quad Y' &= -HL \qquad + FN - BP, \\
Z &= \phantom{-}BL - AM \qquad - HP, &\quad Z' &= \phantom{-}GL - FM \qquad - CP, \\
W &= -FL - GM - HN \qquad, &\quad W' &= -AL - BM - CN \qquad,
\end{align*}
et de l\`{a}, identiquement:
\begin{align*}
\Theta(\Pi^{2} - s^{2})\xi &= \Theta X - \Theta' X', \\
\Theta(\Pi^{2} - s^{2})\eta &= \Theta Y - \Theta' Y', \\
\Theta(\Pi^{2} - s^{2})\zeta &= \Theta Z - \Theta' Z', \\
\Theta(\Pi^{2} - s^{2})\omega &= \Theta W - \Theta' W':
\end{align*}
ce qui correspond en effet au th\'{e}or\`{e}me \'{e}nonc\'{e} ci-dessus, savoir que la r\'{e}ciproque gauche d'un point rencontre les droites $AC$, $BD$ en deux points tels que la droite passant par ces deux points passe par le point m\^{e}me. Les d\'{e}monstrations des diff\'{e}rents th\'{e}or\`{e}mes d'analyse dont je me sers, n'ont point de difficult\'{e}s.



\newpage
% ============================================================
% PAPER 71
% ============================================================
\section*{71. Note sur les fonctions du second ordre.}
\markboth{Note sur les fonctions du second ordre.}{}
\addcontentsline{toc}{section}{71. Note sur les fonctions du second ordre.}

\begin{flushright}
\textit{[From the Journal f"{u}r die reine und angewandte Mathematik (Crelle), tom.\ XXXVIII.\ (1848),\\
pp.\ 105--106.]}
\end{flushright}

\setcounter{page}{421}

\textsc{Soient} $x$, $y$, \ldots\ des variables dont le nombre est $2n$ ou $2n+1$; repr\'{e}sentons par $\xi$, $\eta$, \ldots\ un nombre \'{e}gal de fonctions lin\'{e}aires des variables $x$, $y$, \ldots\ et soit
\[
U = \xi x + \eta y + \cdots = Ax^{2} + By^{2} + \cdots + 2Hxy + \cdots
\]
et
\[
V = \begin{vmatrix}
\xi, & A, & H, & \cdots \\
\eta, & H, & B, & \cdots \\
\vdots & \vdots & \vdots & \ddots
\end{vmatrix}.
\]
J'ai trouv\'{e} que les deux fonctions $U$ et $V$ peuvent \^{e}tre r\'{e}duites \`{a} la forme
\[
U = \lambda\Omega + \mu\Theta, \qquad V = \lambda'\Omega + \mu'\Theta,
\]
o\`{u} $\lambda$, $\mu$, $\lambda'$, $\mu'$ sont des coefficients constants, $\Omega$ est une fonction du second ordre des $n$ variables (fonctions lin\'{e}aires de $x$, $y$, \ldots), et $\Theta$ est une fonction de $n$ ou de $n+1$ variables (fonctions lin\'{e}aires de $x$, $y$, \ldots); selon que le nombre des variables $x$, $y$, \ldots\ est $2n$ ou $2n+1$.

Par exemple pour trois variables $x$, $y$, $z$, on a
\[
U = \lambda x^{2} + \mu yz, \qquad V = \lambda' x^{2} + \mu' yz;
\]
et les coniques repr\'{e}sent\'{e}es par les \'{e}quations $U=0$, $V=0$ ont entre elles un contact double. Cela se rapporte \`{a} la th\'{e}orie des r\'{e}ciproques dans la g\'{e}om\'{e}trie plane, telle que M.~Pl"{u}cker l'a pr\'{e}sent\'{e}e dans son ``System der analytischen Geometrie'' [4$^{\circ}$, Berlin, 1833].

\newpage

\setcounter{page}{422}

De m\^{e}me pour quatre variables $x$, $y$, $z$, $w$ on a
\[
U = \lambda xy + \mu zw, \qquad V = \lambda' xy + \mu' zw,
\]
et les surfaces repr\'{e}sent\'{e}es par les \'{e}quations $U=0$, $V=0$ se coupent dans quatre droites, ou bien se touchent aux quatre sommets d'un quadrilat\`{e}re gauche. Cela se rapporte \'{e}galement \`{a} la th\'{e}orie des r\'{e}ciproques dans l'espace: th\'{e}orie dont j'ai parl\'{e} dans mon M\'{e}moire ``Sur quelques th\'{e}or\`{e}mes de la g\'{e}om\'{e}trie de position,'' \S~VI, [70].

Il y a \`{a} remarquer que dans le cas o\`{u} les coefficients de $x$, $y$, \ldots\ dans les fonctions $\xi$, $\eta$, \ldots\ forment un syst\`{e}me sym\'{e}trique, c'est-\`{a}-dire, o\`{u}
\[
\xi = Ax + Hy + \cdots, \qquad \eta = Hx + By + \cdots,
\]
on a $\lambda:\mu = \lambda':\mu'$, et de l\`{a} $V = KU$, $K$ \'{e}tant donn\'{e} par l'\'{e}quation
\[
K = \begin{vmatrix}
A, & H, & \cdots \\
H, & B, & \cdots \\
\vdots & \vdots & \ddots
\end{vmatrix},
\]
ce qui est connu.

Remarquons enfin que dans le cas g\'{e}n\'{e}ral o\`{u}
\[
\xi = ax + by + \cdots, \qquad \eta = a'x + b'y + \cdots,
\]
la fonction $V$ ne change pas de valeur en \'{e}crivant au lieu de ces expressions:
\[
\xi = ax + a'y + \cdots, \qquad \eta = bx + b'y + \cdots,
\]
propri\'{e}t\'{e} qui se rapporte aussi \`{a} la th\'{e}orie des r\'{e}ciproques.



\newpage
% ============================================================
% PAPER 72
% ============================================================
\section*{72. Note on the theory of permutations.}
\markboth{Note on the theory of permutations.}{}
\addcontentsline{toc}{section}{72. Note on the theory of permutations.}

\begin{flushright}
\textit{[From the Philosophical Magazine, vol.\ xxxiv.\ (1849), pp.\ 527--529.]}
\end{flushright}

\setcounter{page}{423}

\textsc{It} seems worth inquiring whether the distinction made use of in the theory of determinants, of the permutations of a series of things all of them different, into positive and negative permutations, can be made in the case of a series of things not all of them different. The ordinary rule is well known, viz.\ permutations are considered as positive or negative according as they are derived from the primitive arrangement by an even or an odd number of inversions (that is, interchanges of two things); and it is obvious that this rule fails when two or more of the series of things become identical, since in this case any given permutation can be derived indifferently by means of an even or an odd number of inversions. To state the rule in a different form, it will be convenient to enter into some preliminary explanations.

Consider a series of $n$ things, all of them different, and let $abc\ldots$ be the primitive arrangement; imagine a symbol such as $(xyz)(u)(vw)\ldots$ where $x$, $y$, \&c., are the entire series of $n$ things, and which symbol is to be considered as furnishing a rule by which a permutation is to be derived from the primitive arrangement $abc\ldots$ as follows, viz.\ the $(xyz)$ of the symbol denotes that the letters $x$, $y$, $z$ in the primitive arrangement $abc\ldots$ are to be interchanged $x$ into $y$, $y$ into $z$, $z$ into $x$. The $(u)$ of the symbol denotes that the letter $u$ in the primitive arrangement $abc\ldots$ is to remain unaltered. The $(vw)$ of the symbol denotes that the letters $v$, $w$ in the primitive arrangement are to be interchanged $v$ into $w$ and $w$ into $v$, and so on. It is easily seen that any permutation whatever can be derived (and derived in one manner only) from the primitive arrangement by means of a rule such as is furnished by the symbol in question\footnote{See on this subject Cauchy's ``M\'{e}moire sur les Arrangemens \&c.,'' Exercises d'Analyse et de Physique Math\'{e}matique, t.~III. [1844], p.~151.}; and moreover that the number of inversions requisite in order to obtain the permutation by means of the rule in question, is always the smallest number of inversions by which the permutation can be derived. Let $\alpha$, $\beta$ \ldots\ be the number of letters in the components $(xyz)$, $(u)$, $(vw)$, \&c., $\lambda$ the number of these components. The number of inversions in question is evidently $\alpha-1 + \beta-1 + \&$c., or what comes to the same thing, this number is $(n - \lambda)$. It will be convenient to term this number $\lambda$ the \textit{exponent of irregularity} of the permutation, and then $(n - \lambda)$ may be termed the supplement of the exponent of irregularity.

The rule in the case of a series of things, all of them different, may consequently be stated as follows: ``a permutation is positive or negative according as the supplement of the exponent of irregularity is even or odd.''

Consider now a series of things, not all of them different, and suppose that this is derived from the system of the same number of things $abc\ldots$ all of them originally different, by supposing for instance $a = b = \&$c., $f = g = \&$c. A given permutation of the system of things not all of them different, is of course derivable under the supposition in question from several different permutations of the series $abc\ldots$. Considering the supplements of the exponents of irregularity of these last-mentioned several permutations, we may consider the given permutation as positive or negative according as the least of these numbers is even or odd. Hence we obtain the rule, ``a permutation of a series of things not all of them different, is positive or negative according as the minimum supplement of irregularity of the permutation is even or odd, the system being considered as a particular case of a system of the same number of things all of them different, and the given permutation being successively considered as derived from the different permutations which upon this supposition reduce themselves to the given permutation.''

This only differs from the rule, ``a permutation of a series of things, not all of them different, is positive or negative according as the minimum number of inversions by which it can be obtained is even or odd, the system being considered \&c.,'' inasmuch as the former enunciation is based upon and indicates a direct method of determining the minimum number of inversions requisite in order to obtain a given permutation; but the latter is, in simple cases, of the easier application.

As a very simple example, treated by the former rule, we may consider the permutation $1212$ derived from the primitive arrangement $1122$. Considering this primitive arrangement as a particular case of $abcd$, there are four permutations which, on the suppositions $a = b = 1$, $c = d = 2$, reduce themselves to $1212$, viz.\ $acbd$, $bcad$, $adbc$, $bdac$, which are obtained by means of the respective symbols $(a)(bc)(d)$; $(abc)(d)$; $(a)(bdc)$; $(abdc)$, the supplements of the exponents of irregularity being therefore $1$, $2$, $2$, $3$, or the permutation being negative; in fact it is obviously derivable by means of an inversion of the two mean terms.



\newpage
% ============================================================
% PAPER 73
% ============================================================
\section*{73. Abstract of a memoir by Dr Hesse on the construction of the surface of the second order which passes through nine given points.}
\markboth{Abstract of a memoir by Dr Hesse}{}
\addcontentsline{toc}{section}{73. Abstract of a memoir by Dr Hesse on the construction of the surface of the second order.}

\begin{flushright}
\textit{[From the Cambridge and Dublin Mathematical Journal, vol.\ iv.\ (1849), pp.\ 44--46.]}
\end{flushright}

\setcounter{page}{425}

\textsc{The} construction to be presently given of the surface of the second order which passes through nine given points, is taken from a memoir by Dr Hesse (Crelle, t.~xxiv. [1842], p.~36). It depends upon the following lemma, which is there demonstrated.

\begin{lemma}
The polar plane of a fixed point $P$ with respect to any surface of the second order passing through seven given points, passes through a fixed point $Q$ (which may be termed the harmonic pole of the point $P$ with respect to the system of surfaces of the second order).
\end{lemma}

\begin{problem}
Given the seven points $1$, $2$, $3$, $4$, $5$, $6$, $7$, and a point $P$, to construct the harmonic pole $Q$ of the point $P$ with respect to the system of surfaces of the second order passing through the seven points.
\end{problem}

The required point $Q$ may be considered as the intersection of the polar planes of the point $P$ with respect to any three hyperboloids, each of which passes through the seven given points; any such hyperboloid may be considered as determined by means of three of its generating lines. These considerations lead to the construction following.

1. Connecting the points $1$ and $2$, and also the points $3$ and $4$, by two straight lines, and determining the three lines, each of which passes through one of the points $5$, $6$, $7$, and intersects both of the first-mentioned lines, the three lines so determined are generating lines of a hyperboloid passing through the seven points.

Two other systems of generating lines (belonging to two new hyperboloids) are determined by the like construction, interchanging the points $1$, $2$, $3$, $4$. And by interchanging all the seven points we obtain $105$ systems of generating lines (belonging to as many different hyperboloids, unless some of these hyperboloids are identical).

2. It remains to be shown how the polar plane of the point $P$ with respect to one of the $105$ hyperboloids may be constructed. Drawing through the point $P$ three lines, each of which passes through two of the three given generating lines of the hyperboloid in question, the points of intersection of the lines so determined with the generating lines which they respectively intersect, are points of the hyperboloid. Hence, constructing upon each of the three lines in question the harmonic pole of the point $P$ with respect to the two points of intersection, the plane passing through the three harmonic poles is the polar plane of $P$ with respect to the hyperboloid. Hence, constructing the polar planes of $P$ with respect to any three of the $105$ hyperboloids, the point of intersection of these three polar planes is the required point $Q$.

\begin{problem}
To construct the polar plane of a point $P$ with respect to the surface of the second order which passes through nine given points $1$, $2$, $3$, $4$, $5$, $6$, $7$, $8$, $9$.
\end{problem}

Consider any seven of the nine points, e.g.\ the points $1$, $2$, $3$, $4$, $5$, $6$, $7$, and construct the harmonic pole of the point $P$ with respect to the system of surfaces of the second order passing through these seven points. By permuting the different points we obtain $36$ different points $Q$, all of which lie in the same plane. This plane (which is of course determined by any three of the thirty-six points) is the required polar plane. Hence we obtain the solution of

\begin{problem}
To construct the surface of the second order which passes through nine given points $1$, $2$, $3$, $4$, $5$, $6$, $7$, $8$, $9$.
\end{problem}

Assuming the point $P$ arbitrarily, construct the polar plane of this point with respect to the surface of the second order passing through the nine points. Join the point $P$ with any one of the nine points, e.g.\ the point $1$, and on the line so formed determine the harmonic pole $R$ of the point $1$ with respect to the point $P$, and the point where the line $P1$ is intersected by the polar plane. $R$ is a point of the required surface of the second order, which surface is therefore determined by giving every possible position to the point $P$.

This construction is the complete analogue of Pascal's theorem considered as a construction for describing the conic section which passes through five given points. And it would appear that the principles by means of which the construction is obtained ought to lead to the analogue of Pascal's theorem considered in its ordinary form, that is, as a relation between six points of a conic, or in other words to the solution of the problem to determine the relation between ten points of a surface of the second order; but this problem, one of the most interesting in the theory of surfaces of the second order, remains as yet unsolved. The problem last mentioned was proposed as a prize question by the Brussels Academy, which subsequently proposed the more general

\newpage

\setcounter{page}{427}

\begin{center}
\textbf{Construction of the surface of the second order \&c.}
\end{center}

Question to determine the analogue of Pascal's theorem for surfaces of the second order. This of course admitted of being answered in a variety of different ways, according to the different ways of viewing the theorem of Pascal. Thus, M.~Chasles, considering Pascal's theorem as a property of a conic intersected by the three sides of a triangle, discovered the following very elegant analogous theorem for surfaces of the second order.

``The six edges of a tetrahedron may be considered as intersecting a surface of the second order in twelve points lying three and three upon four planes, each one of which contains three points lying on edges which pass through the same angle of the tetrahedron; these planes meet the faces opposite to these angles in four straight lines which are generating lines (of the same species) of a certain hyperboloid.''

It is hardly necessary to remark that all the properties involved in the present memoir are such as to admit of being transformed by the theory of reciprocal polars.



\newpage
% ============================================================
% PAPER 74
% ============================================================
\section*{74. On the simultaneous transformation of two homogeneous functions of the second order.}
\markboth{On the simultaneous transformation of two homogeneous functions of the second order.}{}
\addcontentsline{toc}{section}{74. On the simultaneous transformation of two homogeneous functions of the second order.}

\begin{flushright}
\textit{[From the Cambridge and Dublin Mathematical Journal, vol.\ iv.\ (1849), pp.\ 47--50.]}
\end{flushright}

\setcounter{page}{428}

\textsc{The} theory of the simultaneous transformation by linear substitutions of two homogeneous functions of the second order has been developed by Jacobi in the memoir ``De binis quibuslibet functionibus \&c., Crelle, t.~xii. [1834], p.~1; but the simplest method of treating the problem is the one derived from Mr Boole's Theory of Linear Transformations, combined with the remark in his ``Notes on Linear Transformations,'' in the Cambridge Mathematical Journal, vol.~iv. [1845], p.~167. As I shall have occasion to refer to the results of this theory in the second part of my paper ``On the Attraction of Ellipsoids,'' in the present number of the Journal [75], I take this opportunity of developing the formul\ae\ in question; considering for greater convenience the case of three variables only.

Suppose that by a linear transformation,
\begin{align*}
x &= \alpha x_{1} + \beta y_{1} + \gamma z_{1}, \\
y &= \alpha' x_{1} + \beta' y_{1} + \gamma' z_{1}, \\
z &= \alpha'' x_{1} + \beta'' y_{1} + \gamma'' z_{1}.
\end{align*}
We have identically,
\begin{align*}
ax^{2} + by^{2} + cz^{2} + 2fyz + 2gzx + 2hxy &= a_{1}x_{1}^{2} + b_{1}y_{1}^{2} + c_{1}z_{1}^{2} + 2f_{1}y_{1}z_{1} + 2g_{1}z_{1}x_{1} + 2h_{1}x_{1}y_{1}, \\
Ax^{2} + By^{2} + Cz^{2} + 2Fyz + 2Gzx + 2Hxy &= A_{1}x_{1}^{2} + B_{1}y_{1}^{2} + C_{1}z_{1}^{2} + 2F_{1}y_{1}z_{1} + 2G_{1}z_{1}x_{1} + 2H_{1}x_{1}y_{1}.
\end{align*}
Of course, whatever be the values of $a$, $b$, $c$, $f$, $g$, $h$, the same transformation gives
\[
\alpha(ax^{2} + by^{2} + cz^{2} + \cdots) = \alpha(a_{1}x_{1}^{2} + b_{1}y_{1}^{2} + c_{1}z_{1}^{2} + \cdots),
\]
provided that we have
\begin{align*}
a_{1} &= a\alpha^{2} + b\alpha'^{2} + c\alpha''^{2} + 2f\alpha'\alpha'' + 2g\alpha''\alpha + 2h\alpha\alpha', \\
b_{1} &= a\beta^{2} + b\beta'^{2} + c\beta''^{2} + 2f\beta'\beta'' + 2g\beta''\beta + 2h\beta\beta', \\
c_{1} &= a\gamma^{2} + b\gamma'^{2} + c\gamma''^{2} + 2f\gamma'\gamma'' + 2g\gamma''\gamma + 2h\gamma\gamma', \\
f_{1} &= a\beta\gamma + b\beta'\gamma' + c\beta''\gamma'' + f(\beta'\gamma'' + \beta''\gamma') + g(\beta''\gamma + \beta\gamma'') + h(\beta\gamma' + \beta'\gamma), \\
g_{1} &= a\gamma\alpha + b\gamma'\alpha' + c\gamma''\alpha'' + f(\gamma'\alpha'' + \gamma''\alpha') + g(\gamma''\alpha + \gamma\alpha'') + h(\gamma\alpha' + \gamma'\alpha), \\
h_{1} &= a\alpha\beta + b\alpha'\beta' + c\alpha''\beta'' + f(\alpha'\beta'' + \alpha''\beta') + g(\alpha''\beta + \alpha\beta'') + h(\alpha\beta' + \alpha'\beta).
\end{align*}

\newpage

\setcounter{page}{429}

Representing for a moment the equations between the pairs of functions of the second order by
\[
u = u_{1}, \qquad U = U_{1}, \qquad V = V_{1},
\]
we have, whatever be the value of $\lambda$,
\[
\lambda u + U + V = \lambda u_{1} + U_{1} + V_{1}.
\]
Hence, if
\[
\begin{vmatrix}
\lambda a + A, & \lambda h + H, & \lambda g + G \\
\lambda h + H, & \lambda b + B, & \lambda f + F \\
\lambda g + G, & \lambda f + F, & \lambda c + C
\end{vmatrix}
= \Pi^{3},
\]
then
\[
\begin{vmatrix}
\lambda a_{1} + A_{1}, & \lambda h_{1} + H_{1}, & \lambda g_{1} + G_{1} \\
\lambda h_{1} + H_{1}, & \lambda b_{1} + B_{1}, & \lambda f_{1} + F_{1} \\
\lambda g_{1} + G_{1}, & \lambda f_{1} + F_{1}, & \lambda c_{1} + C_{1}
\end{vmatrix}
= \Pi^{3}.
\]
Hence, since $a$, $b$, $c$, $f$, $g$, $h$, are arbitrary,
\[
\begin{vmatrix}
\lambda a + A, & \lambda h + H, & \lambda g + G \\
\lambda h + H, & \lambda b + B, & \lambda f + F \\
\lambda g + G, & \lambda f + F, & \lambda c + C
\end{vmatrix}
=
\begin{vmatrix}
\lambda a_{1} + A_{1}, & \lambda h_{1} + H_{1}, & \lambda g_{1} + G_{1} \\
\lambda h_{1} + H_{1}, & \lambda b_{1} + B_{1}, & \lambda f_{1} + F_{1} \\
\lambda g_{1} + G_{1}, & \lambda f_{1} + F_{1}, & \lambda c_{1} + C_{1}
\end{vmatrix},
\]
which determine the relations which must subsist between the coefficients of the functions of the second order. Whence we derive
\[
\begin{vmatrix}
a, & h, & g \\
h, & b, & f \\
g, & f, & c
\end{vmatrix}
= \Pi^{3}
\begin{vmatrix}
a_{1}, & h_{1}, & g_{1} \\
h_{1}, & b_{1}, & f_{1} \\
g_{1}, & f_{1}, & c_{1}
\end{vmatrix},
\]
and by comparing the coefficients of $\alpha$, \&c., if we write for shortness,
\[
\mathfrak{A} = (\lambda b + B)(\lambda c + C) - (\lambda f + F)^{2}, \quad \&\text{c.},
\]
then
\begin{align*}
\mathfrak{A}\alpha^{2} + \mathfrak{B}\beta^{2} + \mathfrak{C}\gamma^{2} + 2\mathfrak{F}\beta\gamma + 2\mathfrak{G}\gamma\alpha + 2\mathfrak{H}\alpha\beta &= \Pi^{3}, \\
\mathfrak{A}\alpha'^{2} + \mathfrak{B}\beta'^{2} + \mathfrak{C}\gamma'^{2} + 2\mathfrak{F}\beta'\gamma' + 2\mathfrak{G}\gamma'\alpha' + 2\mathfrak{H}\alpha'\beta' &= \Pi^{3}, \\
\mathfrak{A}\alpha''^{2} + \mathfrak{B}\beta''^{2} + \mathfrak{C}\gamma''^{2} + 2\mathfrak{F}\beta''\gamma'' + 2\mathfrak{G}\gamma''\alpha'' + 2\mathfrak{H}\alpha''\beta'' &= \Pi^{3}, \\
\mathfrak{A}\alpha'\alpha'' + \mathfrak{B}\beta'\beta'' + \mathfrak{C}\gamma'\gamma'' + \mathfrak{F}(\beta'\gamma'' + \beta''\gamma') + \mathfrak{G}(\gamma'\alpha'' + \gamma''\alpha') + \mathfrak{H}(\alpha'\beta'' + \alpha''\beta') &= \Pi^{3}\mathfrak{F}_{1}, \\
\mathfrak{A}\alpha''\alpha + \mathfrak{B}\beta''\beta + \mathfrak{C}\gamma''\gamma + \mathfrak{F}(\beta''\gamma + \beta\gamma'') + \mathfrak{G}(\gamma''\alpha + \gamma\alpha'') + \mathfrak{H}(\alpha''\beta + \alpha\beta'') &= \Pi^{3}\mathfrak{G}_{1}, \\
\mathfrak{A}\alpha\alpha' + \mathfrak{B}\beta\beta' + \mathfrak{C}\gamma\gamma' + \mathfrak{F}(\beta\gamma' + \beta'\gamma) + \mathfrak{G}(\gamma\alpha' + \gamma'\alpha) + \mathfrak{H}(\alpha\beta' + \alpha'\beta) &= \Pi^{3}\mathfrak{H}_{1},
\end{align*}
each of which virtually contains three equations on account of the indeterminate quantity $\lambda$. A somewhat more elegant form may be given to these equations; thus the first of them is
\[
\begin{vmatrix}
\mathfrak{A}, & \lambda a_{1} + A_{1}, & \lambda h_{1} + H_{1}, & \lambda g_{1} + G_{1} \\
\alpha, & \lambda a + A, & \lambda h + H, & \lambda g + G \\
\beta, & \lambda h + H, & \lambda b + B, & \lambda f + F \\
\gamma, & \lambda g + G, & \lambda f + F, & \lambda c + C
\end{vmatrix}
= \Pi^{3},
\]
from which the form of the whole system is sufficiently obvious. The actual values of the coefficients $\alpha$, $\beta$, \&c.\ can only be obtained in the particular case where
\[
a_{1} = b_{1} = c_{1} = 1, \qquad f_{1} = g_{1} = h_{1} = 0,
\]
or (what is the same thing) where
\[
A_{1} = a, \quad B_{1} = b, \quad C_{1} = c, \quad F_{1} = f, \quad G_{1} = g, \quad H_{1} = h.
\]
The preceding formul\ae\ are, however, simplified by assuming as the additional condition
\[
(\lambda a + A)(\lambda b + B)(\lambda c + C) = \Pi^{3},
\]
or $\Pi^{3} = \Theta^{-1}$ suppose; and then
\begin{align*}
\mathfrak{A} &= (\lambda b + B)(\lambda c + C) - (\lambda f + F)^{2} = \mathfrak{a}, \\
\mathfrak{B} &= (\lambda c + C)(\lambda a + A) - (\lambda g + G)^{2} = \mathfrak{b}, \\
\mathfrak{C} &= (\lambda a + A)(\lambda b + B) - (\lambda h + H)^{2} = \mathfrak{c}, \\
\mathfrak{F} &= (\lambda g + G)(\lambda h + H) - (\lambda a + A)(\lambda f + F) = \mathfrak{f}, \\
\mathfrak{G} &= (\lambda h + H)(\lambda f + F) - (\lambda b + B)(\lambda g + G) = \mathfrak{g}, \\
\mathfrak{H} &= (\lambda f + F)(\lambda g + G) - (\lambda c + C)(\lambda h + H) = \mathfrak{h}.
\end{align*}

\newpage

\setcounter{page}{431}

where, writing down the expanded values of $\mathfrak{a}$, $\mathfrak{b}$, $\mathfrak{c}$, $\mathfrak{f}$, $\mathfrak{g}$, $\mathfrak{h}$,
\begin{align*}
(\lambda b + B)(\lambda c + C) - (\lambda f + F)^{2} &= \mathfrak{a}, \\
(\lambda c + C)(\lambda a + A) - (\lambda g + G)^{2} &= \mathfrak{b}, \\
(\lambda a + A)(\lambda b + B) - (\lambda h + H)^{2} &= \mathfrak{c}, \\
(\lambda g + G)(\lambda h + H) - (\lambda a + A)(\lambda f + F) &= \mathfrak{f}, \\
(\lambda h + H)(\lambda f + F) - (\lambda b + B)(\lambda g + G) &= \mathfrak{g}, \\
(\lambda f + F)(\lambda g + G) - (\lambda c + C)(\lambda h + H) &= \mathfrak{h}.
\end{align*}
By writing successively $\lambda = -A_{1}$, $\lambda = -B_{1}$, $\lambda = -C_{1}$, we see in the first place that $A_{1}$, $B_{1}$, $C_{1}$ are the roots of the same cubic equation, and we obtain next the values of $\alpha^{2}$, $\beta^{2}$, $\gamma^{2}$, \&c.\ in terms of these quantities $A_{1}$, $B_{1}$, $C_{1}$, and of the coefficients $a$, $b$, \&c., $A$, $B$, \&c. It is easy to see how the above formul\ae\ would have been modified if $a_{1}$, $b_{1}$, $c_{1}$, instead of being equal to unity, had one or more of them been equal to unity with a negative sign. It is obvious that every step of the preceding process is equally applicable whatever be the number of variables.



\newpage
% ============================================================
% PAPER 75
% ============================================================
\section*{75. On the attraction of an ellipsoid.}
\markboth{On the attraction of an ellipsoid.}{}
\addcontentsline{toc}{section}{75. On the attraction of an ellipsoid.}

\begin{flushright}
\textit{[From the Cambridge and Dublin Mathematical Journal, vol.\ iv.\ (1849), pp.\ 50--65.]}
\end{flushright}

\setcounter{page}{432}

\begin{center}
\textbf{Part I.---On Legendre's Solution of the Problem of the Attraction of an Ellipsoid on an External Point.}
\end{center}

\textsc{I} propose in the following paper to give an outline of Legendre's investigation of the attraction of an ellipsoid upon an exterior point, [``M\'{e}moire sur les Int\'{e}grales Doubles,'' Paris, M\'{e}m.\ Acad.\ Sc.\ for 1788, published 1791, pp.~454--486], one of the earliest and (notwithstanding its complexity) most elegant solutions of the problem. It will be convenient to begin by considering some of the geometrical properties of a system of cones made use of in the investigation.

\subsection*{\S~1.}

The equation of the ellipsoid referred to axes parallel to the principal axes, and passing through the attracted point, may be written under the form
\[
l(x-a)^{2} + m(y-b)^{2} + n(z-c)^{2} - k = 0,
\]
(where $\frac{k}{l}$, $\frac{k}{m}$, $\frac{k}{n}$ are the semiaxes, and $a$, $b$, $c$ are the coordinates of the attracted point referred to the principal axes). Or putting $la^{2} + mb^{2} + nc^{2} - k = \delta$, this equation becomes
\[
lx^{2} + my^{2} + nz^{2} - 2(lax + mby + ncz) + \delta = 0.
\]

The cones in question are those which have the same axes and directions of circular section as the cone having its vertex in the attracted point and circumscribed about the ellipsoid. The equation of the system of cones (containing the arbitrary parameter $\omega$) is
\[
(lx^{2} + my^{2} + nz^{2})k - (lax + mby + ncz)^{2} + \omega^{2}(x^{2} + y^{2} + z^{2}) = 0;
\]
or as it may also be written,
\[
(\omega^{2} + lk - l^{2}a^{2})x^{2} + (\omega^{2} + mk - m^{2}b^{2})y^{2} + (\omega^{2} + nk - n^{2}c^{2})z^{2} - 2mnbc\,yz - 2nlca\,zx - 2lmab\,xy = 0.
\]

\newpage

\setcounter{page}{433}

For $\omega = 0$, the cone coincides with the circumscribed cone; as $\omega$ increases, the aperture of the cone gradually diminishes, until for a certain value, $\omega = \Omega$, the cone reduces itself to a straight line (the normal of the confocal ellipsoid through the attracted point). It is easily seen that $\Omega^{2}$ is the positive root of the equation
\[
\frac{l^{2}a^{2}}{\omega^{2} + lk} + \frac{m^{2}b^{2}}{\omega^{2} + mk} + \frac{n^{2}c^{2}}{\omega^{2} + nk} = 1,
\]
a different form of which may be obtained by writing $\omega^{2} = \frac{k}{\theta}$, $\theta$ being then determined by means of the equation
\[
\frac{la^{2}}{k + l\theta} + \frac{mb^{2}}{k + m\theta} + \frac{nc^{2}}{k + n\theta} = 1,
\]
$\sqrt{k + l\theta}$, $\sqrt{k + m\theta}$, $\sqrt{k + n\theta}$ are the semiaxes of the confocal ellipsoid through the attracted point.

In the case where $\omega$ remains indeterminate, it is obvious that the cone intersects the ellipsoid in the curve in which the ellipsoid is intersected by a certain hyperboloid of revolution of two sheets, having the attracted point for a focus, and the plane of contact of the ellipsoid with the circumscribed cone (that is the polar plane of the attracted point) for the corresponding directrix plane: also the excentricity of the hyperboloid is $\frac{1}{\omega}\sqrt{l^{2}a^{2} + m^{2}b^{2} + n^{2}c^{2}}$, which suffices for its complete determination. For $\omega = 0$, the hyperboloid reduces itself to the plane of contact of the ellipsoid with the circumscribed cone, and for $\omega = \Omega$, the hyperboloid and the ellipsoid have a double contact, viz.\ at the points where the ellipsoid is intersected by the normal to the confocal ellipsoid through the attracted point.

If $\omega$ remains constant while $k$ is supposed to vary, that is, if the ellipsoid vary in magnitude (the position and proportion of its axes remaining unaltered), the locus of the intersection of the cone and the ellipsoid is a surface of the fourth order defined by the equation
\[
(lx^{2} + my^{2} + nz^{2} - lax - mby - ncz)^{2} = \omega^{2}(x^{2} + y^{2} + z^{2}),
\]
and consisting of an exterior and an interior sheet meeting at the attracted point, which is a conical point on the surface, viz.\ a point where the tangent plane is replaced by a tangent cone. The general form of this surface is easily seen from the figure, in which the ellipsoid has been replaced by a sphere, and the surface in question is that generated by the revolution of the curve round the line $CM$.

\newpage

\setcounter{page}{434}

The surface of the fourth order being once described for any particular value of $\omega$, the cone corresponding to any one of the series of similar, similarly situated, and concentric ellipsoids is at once determined by means of the intersection of the ellipsoid in question with the surface of the fourth order. It is clear too that there is always one of these ellipsoids which has a double contact with the surface of the fourth order, viz.\ at the points where this ellipsoid is intersected by the normal to the confocal ellipsoid through the attracted point; thus there is always an ellipsoid for which the cone corresponding to a given value of $\omega$ reduces itself to a straight line.

Consider the attracting ellipsoid, which for distinction may be termed the ellipsoid $S$, and the two cones $C$, $C'$, which correspond to the values $\omega$, $\omega - d\omega$ of the variable parameter. Legendre shows that the attraction of the portion of the ellipsoid $S$ included between the two cones $C$, $C'$ is independent of the quantity $k$, which determines the magnitude of the ellipsoid: that is, if there be any other ellipsoid $T$ similarly situated and concentric to and with the ellipsoid $S$, and two cones $D$, $D'$, which for the ellipsoid $T$ correspond to the same values $\omega$, $\omega - d\omega$ of the variable parameter; then the attraction of the portion of the ellipsoid $S$, included between the two cones $C$, $C'$, is equal to the attraction of the portion of the ellipsoid $T$ included between the two cones $D$ and $D'$. By taking for the ellipsoid $T$ the ellipsoid for which the cone $D$ reduces itself to a straight line, the aperture of the cone $D'$ is indefinitely small, and the attraction of the portion of the ellipsoid $T$ included within the cone $D'$ is at once determined; and thus the attraction of the portion of the ellipsoid $S$ included between the cones $C$, $C'$ is obtained in a finite form. Hence the attraction of the portion of the ellipsoid $S$ included between any two cones $C_{1}$, $C_{2}$, corresponding to the values $\omega_{1}$, $\omega_{2}$, of the variable parameter, is expressed by means of a single integral, and by extending the integration from $\omega = 0$ to $\omega = \sqrt{\frac{k}{\theta_{1}}}$, the attraction of the whole ellipsoid is obtained in the form of a single integral readily reducible to that given by the ordinary solutions. It is clear too that the attraction of the portion of the ellipsoid $S$ included between any two cones $C_{1}$, $C_{2}$, is equal to that of the portion of the ellipsoid $T$ included between the corresponding cones $D_{1}$ and $D_{2}$. Hence also, assuming for the ellipsoid $T$, that for which the cone $D_{2}$ reduces itself to a straight line, and supposing that the cones $C_{1}$ and $D_{1}$ coincide with the circumscribing cones, the attraction of the portion of the ellipsoid $S$ exterior to the cone $C_{1}$ is equal to the attraction of the entire ellipsoid $T$. More generally, the attraction of the portion of the ellipsoid $S$ included between the cones $C_{1}$ and $C_{2}$ is equal to the attraction of the shell included between the surfaces of the two ellipsoids, for which the cones $D_{1}$ and $D_{2}$ respectively reduce themselves to straight lines.

\subsection*{\S~2.}

Proceeding to the analytical solution, and resuming the equation of the ellipsoid
\[
lx^{2} + my^{2} + nz^{2} - 2(lax + mby + ncz) + \delta = 0,
\]
and that of the cone
\[
(lx^{2} + my^{2} + nz^{2})k - (lax + mby + ncz)^{2} + \omega^{2}(x^{2} + y^{2} + z^{2}) = 0;
\]

\newpage

\setcounter{page}{435}

consider a radius vector on the conical surface such that the cosines of its inclinations to the axes are
\[
\frac{P}{\Theta}, \quad \frac{Q}{\Theta}, \quad \frac{R}{\Theta}, \quad \text{where } \Theta = \sqrt{P^{2} + Q^{2} + R^{2}},
\]
$P$, $Q$, $R$ and $\Theta$ being functions of the parameter $\omega$, and of another variable $\phi$, which determines the position of the radius vector upon the conical surface. Also let $\rho$ be the length of the portion of the radius vector which lies within the ellipsoid; then representing by $dS$ the spherical angle corresponding to the variations of $\omega$ and $\phi$, the attraction in the direction of the axis of $x$ is given by the formula
\[
A = \iint \rho^{3} \frac{P}{\Theta}\, dS.
\]
Also by a known formula
\[
dS = \frac{1}{\Theta^{4}} \left\{ \Theta^{2}\left(\frac{dR}{d\phi}\frac{dQ}{d\omega} - \frac{dR}{d\omega}\frac{dQ}{d\phi}\right) + \cdots \right\} d\omega\, d\phi,
\]
and it is easy to obtain
\[
2\omega\Theta^{4} = lP^{2} + mQ^{2} + nR^{2}.
\]
The quantities $P$, $Q$, $R$ have now to be expressed as functions of $\omega$, $\phi$, so that their values substituted for $x$, $y$, $z$, may satisfy identically the equation of the cone. This may be done by assuming
\begin{align*}
P &= \sqrt{(\omega^{2} + m\delta)(\omega^{2} + n\delta) - m^{2}n^{2}b^{2}c^{2}}, \\
Q &= mb(\omega^{2} + n\delta)\left(\sqrt{\frac{n}{m}}\frac{ca}{\rho'}\cos\phi\right) - nc\sqrt{(\omega^{2} + m\delta)}\sin\phi, \\
R &= nc(\omega^{2} + m\delta)\left(\sqrt{\frac{m}{n}}\frac{ab}{\rho'}\cos\phi\right) - mb\sqrt{(\omega^{2} + n\delta)}\sin\phi,
\end{align*}
where
\begin{align*}
\rho'^{2} &= m^{2}b^{2}(\omega^{2} + n\delta) + n^{2}c^{2}(\omega^{2} + m\delta), \\
D' &= (\omega^{2} + l\delta)(\omega^{2} + m\delta)(\omega^{2} + n\delta),
\end{align*}
a system of values which, in point of fact, depend upon the following geometrical considerations: by treating $x$ as a constant in the equation of the cone, that is, in effect by considering the sections of the cone by planes parallel to that of $yz$, the equation of the cone becomes that of an ellipse; transforming first to a set of axes through the centre and then to a set of conjugate axes, one of which passes through the point where the plane of the ellipse is intersected by the axis of $x$, then the equation takes the form $\frac{\xi^{2}}{A^{2}} + \frac{\eta^{2}}{B^{2}} = 1$, and is satisfied by $\xi = A\cos\phi$, $\eta = B\sin\phi$, \&c., and of course linear functions of these values, the preceding expressions may be obtained.

\newpage

\setcounter{page}{436}

The substitution of the above values of $P$, $Q$, $R$ (a somewhat tedious one which does not occur in the process actually made use of by Legendre) gives the very simple result,
\[
dS = \frac{P\, d\omega\, d\phi}{\omega(lP^{2} + mQ^{2} + nR^{2})};
\]
and the formula for the attraction becomes
\[
A = 2\iint \frac{P^{4}\, d\omega\, d\phi}{lP^{2} + mQ^{2} + nR^{2}},
\]
which is of the form $A = 2\int I\, d\omega$, where
\[
I = \int \frac{P^{4}\, d\phi}{lP^{2} + mQ^{2} + nR^{2}},
\]
which last integral, taken between the limits $\phi = 0$ and $\phi = 2\pi$, and multiplied by $2\omega\, d\omega$, expresses the attraction of the portion of the ellipsoid included between two consecutive cones. The integration is evidently possible, but the actual performance of it is the great difficulty of Legendre's process. The result, as before mentioned, is independent of the quantity $k$, or, what comes to the same thing, of the quantity $\delta$: assuming this property (an assumption which in fact resolves itself into the consideration of the ellipsoid for which the cone reduces itself to a straight line, as before explained), the integral is at once obtained by writing $\delta = A$ where $A$ represents the positive root of the equation
\[
\frac{la^{2}}{\omega^{2} + lA} + \frac{mb^{2}}{\omega^{2} + mA} + \frac{nc^{2}}{\omega^{2} + nA} = 1.
\]
This gives
\begin{align*}
P &= \sqrt{la^{2}(\omega^{2} + mA)(\omega^{2} + nA)}, \\
Q &= l^{2}mab(\omega^{2} + nA), \\
R &= l^{2}nac(\omega^{2} + mA),
\end{align*}
values independent of $\phi$, or the value of $I$ is found by multiplying the quantity under the integral sign by $2\pi$; and hence we have
\[
A = 4\pi la \int_{0}^{\omega_{1}} \frac{(\omega^{2} + lA)^{\frac{1}{2}}(\omega^{2} + mA)^{\frac{1}{2}}(\omega^{2} + nA)^{\frac{1}{2}}\, d\omega}{l^{2}a^{2}(\omega^{2} + mA)^{2}(\omega^{2} + nA)^{2} + m^{2}b^{2}(\omega^{2} + nA)^{2}(\omega^{2} + lA)^{2} + n^{2}c^{2}(\omega^{2} + lA)^{2}(\omega^{2} + mA)^{2}},
\]
where of course $A$ is to be considered as a function of $\omega$. By integrating from $\omega = 0$, to $\omega = \omega_{1}$, we have the attraction of the portion of the ellipsoid included between any two of the series of cones, and to obtain the attraction of the whole ellipsoid we must integrate from $\omega = 0$ to $\omega = \sqrt{\frac{k}{\theta_{1}}}$, where $\theta_{1}$ is determined as before by the equation
\[
\frac{la^{2}}{k + l\theta} + \frac{mb^{2}}{k + m\theta} + \frac{nc^{2}}{k + n\theta} = 1.
\]

\newpage

\setcounter{page}{437}

and it is obvious that for this value of $\omega$ we have $A = \delta$. The expression for the attraction is easily reduced to a known form by writing $y = \frac{k}{\omega^{2}}$; this gives
\[
A = 4\pi la \int \frac{(k + lA)^{\frac{1}{2}}(k + mA)^{\frac{1}{2}}(k + nA)^{\frac{1}{2}}\, dy}{(k + mA)^{2}(k + nA)^{2} + m^{2}b^{2}(k + nA)^{2}(k + lA)^{2} + n^{2}c^{2}(k + lA)^{2}(k + mA)^{2}}.
\]
Also
\[
A = \frac{1}{y}\left(\frac{l^{2}a^{2}}{k + ly} + \frac{m^{2}b^{2}}{k + my} + \frac{n^{2}c^{2}}{k + ny}\right)^{-1},
\]
whence
\[
\frac{dA}{dy} = \frac{2(k + ly)^{-2}(k + my)^{-2}(k + ny)^{-2}}{l^{2}a^{2}(k + my)^{2}(k + ny)^{2} + \cdots},
\]
and thus
\[
A = 2\pi l^{\frac{1}{2}}a \int \frac{dy}{(k + ly)^{\frac{1}{2}}(k + my)^{\frac{1}{2}}(k + ny)^{\frac{1}{2}}},
\]
where for the entire ellipsoid the integral is to be taken from $y = \theta_{1}$ to $y = \infty$. A better known form is readily obtained by writing $y = kr - \frac{1}{l}\frac{1}{1 + x^{2}}$, in which case the limits for the entire ellipsoid are $x = 0$, $x = 1$.

It may be as well to indicate the first step of the reduction of the integral $I$, viz.\ the method of resolving the denominator into two factors. We have identically,
\[
(A - \delta)(lP^{2} + mQ^{2} + nR^{2}) = \omega^{2}(P^{2} + Q^{2} + R^{2}) + A(lP^{2} + mQ^{2} + nR^{2}) - (laP + mbQ + ncR)^{2},
\]
and the second side of this equation is resolvable into two factors independently of the particular values of $P$, $Q$, $R$. Representing this second side for a moment in the notation of a general quadratic function, or under the form
\[
AP^{2} + BQ^{2} + CR^{2} + 2FQR + 2GRP + 2HPQ,
\]
we have the required solution,
\begin{multline*}
lP^{2} + mQ^{2} + nR^{2} = \frac{1}{A}\left[\mathfrak{A}P + \{\mathfrak{F} + \sqrt{(-\mathfrak{B})}\}Q + \{\mathfrak{G} + \sqrt{(-\mathfrak{C})}\}R\right] \\
\times \left[\mathfrak{A}P + \{\mathfrak{F} - \sqrt{(-\mathfrak{B})}\}Q + \{\mathfrak{G} - \sqrt{(-\mathfrak{C})}\}R\right];
\end{multline*}
where, as usual, $\mathfrak{B} = CA - G^{2}$, $\mathfrak{C} = AB - H^{2}$, and the roots must be so taken that $\sqrt{(-\mathfrak{B})}\sqrt{(-\mathfrak{C})} = \mathfrak{F}$ ($\mathfrak{F} = GH - A\mathfrak{F}$).

I have purposely restricted myself so far to the problem considered by Legendre: the general transformation, of which the preceding is a particular case, and also a simpler mode of effecting the integration, are given in the next part of this paper.



\newpage
% ============================================================
\begin{center}
\textbf{Part II.---On a Formula for the Transformation of Certain Multiple Integrals.}
\end{center}

\setcounter{page}{438}

Consider the integral
\[
V = \int F(x, y, \ldots)\, dx\, dy\ldots,
\]
where the number of variables $x$, $y$, \ldots\ is equal to $n$, and $F(x, y, \ldots)$ is a homogeneous function of the order $\mu$.

Suppose that $x$, $y$, \ldots\ are connected by a homogeneous equation $\psi(x, y, \ldots) = 0$ containing a variable parameter $\omega$ (so that $\omega$ is a homogeneous function of the order zero in the variables $x$, $y$, \ldots). Then, writing
\[
r^{2} = x^{2} + y^{2} + \ldots, \qquad x = r\alpha, \quad y = r\beta, \ldots
\]
the quantities $\alpha$, $\beta$, \ldots\ are connected by the equations
\[
\alpha^{2} + \beta^{2} + \ldots = 1, \qquad \psi(\alpha, \beta, \ldots) = 0,
\]
and we may therefore consider them as functions of $\omega$ and of $(n-2)$ independent variables $\theta$, $\phi$, \&c.; whence
\[
dx\, dy\ldots = r^{n-1}\, V\, dr\, d\omega\, d\theta\ldots,
\]
where
\[
V = \begin{vmatrix}
\alpha, & \beta, & \ldots \\
\frac{d\alpha}{d\omega}, & \frac{d\beta}{d\omega}, & \ldots \\
\frac{d\alpha}{d\theta}, & \frac{d\beta}{d\theta}, & \ldots \\
\vdots & \vdots & \ddots
\end{vmatrix}.
\]
Also $F(x, y, \ldots) = r^{\mu}F(\alpha, \beta, \ldots)$, and therefore
\[
V = \int r^{\mu+n-1}\, F(\alpha, \beta, \ldots)\, V\, dr\, d\omega\, d\theta\ldots,
\]
or, integrating with respect to $r$,
\[
\int r^{\mu+n-1}\, dr = \frac{r^{\mu+n}}{\mu+n},
\]
which, taken between the proper limits, is a function of $\alpha$, $\beta$, \ldots, equal $f(\alpha, \beta, \ldots)$ suppose; this gives
\[
V = \int f(\alpha, \beta, \ldots)F(\alpha, \beta, \ldots)\, V\, d\omega\, d\theta\ldots,
\]
in which I shall assume that the limits of $\omega$ are constant. If, in order to get rid of the condition $\alpha^{2} + \beta^{2} + \ldots = 1$, we assume

\newpage

\setcounter{page}{439}

the preceding expression for $V$ becomes
\[
V = \int f\left(\frac{p}{\sqrt{p^{2} + q^{2} + \ldots}}, \frac{q}{\sqrt{p^{2} + q^{2} + \ldots}}, \ldots\right) F(p, q, \ldots)\, \frac{D}{\Delta}\, d\omega\, d\theta\ldots,
\]
where
\[
D = \begin{vmatrix}
p, & q, & \ldots \\
\frac{dp}{d\omega}, & \frac{dq}{d\omega}, & \ldots \\
\frac{dp}{d\theta}, & \frac{dq}{d\theta}, & \ldots \\
\vdots & \vdots & \ddots
\end{vmatrix}, \qquad
\Delta = \sqrt{p^{2} + q^{2} + \ldots}.
\]
Assume
\begin{align*}
p &= P\xi + P'\eta + P''\zeta + \ldots, \\
q &= Q\xi + Q'\eta + Q''\zeta + \ldots,
\end{align*}
where the number of variables $\xi$, $\eta$, $\zeta$, \ldots\ (functions in general of $\omega$, $\theta$, \&c.) is $\mu$, and where the coefficients $P$, $Q$, \&c.\ are supposed to be functions of $\omega$ only. We have
\[
\frac{dp}{d\theta} = P\frac{d\xi}{d\theta} + P'\frac{d\eta}{d\theta} + \ldots,
\]
and, substituting these values as well as those of $p$, $q$, \&c., but retaining the terms $\frac{dp}{d\omega}$, $\frac{dq}{d\omega}$, \&c., in their original form, the determinant $D$ resolves itself into the sum of a series of products,
\[
D = \sum \begin{vmatrix}
\frac{dp}{d\omega}, & \frac{dq}{d\omega}, & \ldots \\
P', & Q', & \ldots \\
P'', & Q'', & \ldots \\
\vdots & \vdots & \ddots
\end{vmatrix}
\begin{vmatrix}
\xi, & \eta, & \ldots \\
\frac{d\xi}{d\theta}, & \frac{d\eta}{d\theta}, & \ldots \\
\vdots & \vdots & \ddots
\end{vmatrix}.
\]
Let $\Psi$ be the function to which $\psi(p, q, \ldots)$ is changed by the substitution of the above values of $p$, $q$, \ldots\ so that $\Psi$ is a homogeneous function of $\xi$, $\eta$, $\zeta$, \ldots\ and we have the relation $\Psi = 0$. (It will be convenient to consider $\xi$, $\eta$, $\zeta$, \ldots\ as functions of $\omega$, $\theta$, \&c., such as to satisfy identically this last equation.) We deduce
\[
\frac{d\xi}{d\theta} = \frac{d\Psi}{d\eta} \div \frac{d\Psi}{d\xi}, \quad \&\text{c.} = \mathcal{S} \text{ suppose}.
\]

\newpage

\setcounter{page}{440}

where for shortness $X = \frac{d\Psi}{d\xi}$, $Y = \frac{d\Psi}{d\eta}$, \&c. The substitution of these values gives
\[
D = \begin{vmatrix}
\frac{dp}{d\omega}, & \frac{dq}{d\omega}, & \ldots \\
X, & P, & P', & \ldots \\
Y, & Q, & Q', & \ldots \\
\vdots & \vdots & \vdots & \ddots
\end{vmatrix},
\]
where it will be remarked that the successive horizontal lines (after the first) of the determinant are the differential coefficients of $\Psi$, $p$, $q$, \ldots\ with respect to $\xi$, with respect to $\eta$, \&c.

In general, if $\psi$ denote any function of $p$, $q$, \ldots\ these quantities being themselves functions of $\omega$, $\xi$, $\eta$, \&c., and $\xi$, $\eta$, \ldots\ containing $\omega$; also if $\Psi$ be what $\psi$ becomes when for $p$, $q$, \ldots\ we substitute their values in $\omega$, $\xi$, $\eta$, \ldots\; then we have identically
\[
\begin{vmatrix}
\frac{dp}{d\omega}, & \frac{dq}{d\omega}, & \ldots \\
X, & P, & P', & \ldots \\
Y, & Q, & Q', & \ldots \\
\vdots & \vdots & \vdots & \ddots
\end{vmatrix}
= \frac{d\Psi}{d\omega}\begin{vmatrix}
P, & P', & \ldots \\
Q, & Q', & \ldots \\
\vdots & \vdots & \ddots
\end{vmatrix}.
\]
In other words, the preceding equation is satisfied by substituting $\frac{d\Psi}{d\omega}$ for $\frac{dp}{d\omega}$, $\frac{dq}{d\omega}$, \ldots\ on the second side, $\frac{d}{d\omega}$ denoting a partial differential coefficient taken only so far as $\omega$ is explicitly contained in $\Psi$. And considering $p$, $q$, \ldots\ as functions of $\omega$, $\xi$, $\eta$, \ldots\ ($\xi$, $\eta$, \ldots\ themselves functions of $\omega$ and of other variables which need not here be considered), $\Psi$ or $\psi$ vanishes identically, and we have $\frac{d\Psi}{d\omega} = 0$. Hence, in the last formula, we have to write $-\frac{d\psi}{d\omega}$ instead of $\frac{dp}{d\omega}$, \&c., and we thus derive
\[
D = \begin{vmatrix}
-\frac{d\psi}{d\omega}, & \ldots \\
X, & P, & P', & \ldots \\
Y, & Q, & Q', & \ldots \\
\vdots & \vdots & \vdots & \ddots
\end{vmatrix}
\begin{vmatrix}
P, & P', & \ldots \\
Q, & Q', & \ldots \\
\vdots & \vdots & \ddots
\end{vmatrix}^{-1}.
\]
This formula, or one equivalent to it, is given in Jacobi's memoir ``De Determinantibus Functionalibus,'' Crelle, t.~XXII. [1841] p.~319.

\newpage

\setcounter{page}{441}

whence $D$ becomes
\[
D = \begin{vmatrix}
P, & P', & \ldots \\
Q, & Q', & \ldots \\
\vdots & \vdots & \ddots
\end{vmatrix}
\frac{d\omega}{d\Psi},
\]
which is the value to be made use of in the equation
\[
V = \int f(\alpha, \beta, \ldots)F(\alpha, \beta, \ldots)\, V\, d\omega\, d\theta\ldots.
\]
The principal use of the formula is where $\psi$ is a homogeneous function of the second order of $p$, $q$, \ldots. Thus, suppose
\[
\psi(p, q, \ldots) = Ap^{2} + Bq^{2} + \ldots + 2Hpq + \ldots;
\]
also, for the sake of conformity to the usual notation in the theory of transformation of quadratic functions, writing $\alpha$, $\alpha'$, \ldots\ $\beta$, $\beta'$, \ldots\ instead of $P$, $P'$, \ldots\ $Q$, $Q'$, \ldots\ and putting after the differentiations $\xi = 1$, we have
\[
p = \alpha + \alpha'\eta + \alpha''\zeta + \ldots,
\]
values which we may assume to give rise to the equation
\[
(Ap^{2} + Bq^{2} \ldots + 2Hpq \ldots) = (1 - \eta^{2} - \zeta^{2} - \ldots),
\]
(where $\eta$, $\zeta$, \ldots\ are taken to be functions of $\theta$, \&c.\ such as to satisfy identically the equation $1 = \eta^{2} + \zeta^{2} + \ldots$).

Hence, by a well-known property, if
\[
\begin{vmatrix}
A, & H, & \ldots \\
H, & B, & \ldots \\
\vdots & \vdots & \ddots
\end{vmatrix}
= (-)^{n-1}\Delta,
\]
we have
\[
\frac{d\Psi}{d\omega} = \frac{1}{\Delta}\frac{d\Delta}{d\omega}(1 - \eta^{2} - \zeta^{2} - \ldots)^{\frac{n}{2}},
\]
so that, observing that in the present case $X = \frac{d\Psi}{d\xi}$, and therefore
\[
\frac{V}{\Delta} = \frac{1}{\Delta}\frac{d\Delta}{d\omega}(1 - \eta^{2} - \zeta^{2} - \ldots)^{\frac{n}{2}-1},
\]
we have
\[
V = \int \frac{f(\alpha, \beta, \ldots)}{\Delta}\frac{d\Delta}{d\omega}(1 - \eta^{2} - \zeta^{2} - \ldots)^{\frac{n}{2}-1}\, d\omega\, d\theta\ldots.
\]

\newpage

\setcounter{page}{442}

The remainder of the process of integration may in many cases be effected by the method made use of by Jacobi in the memoir ``De binis quibuslibet functionibus \&c.'' Crelle, t.~XII. [1834] p.~1, viz.\ the coefficients $\alpha$, $\alpha'$, \&c., $\beta$, \&c.,\ may in addition to the conditions which they are already supposed to satisfy, be so determined as to reduce any homogeneous function of $p$, $q$, $r$, \ldots\ entering into the integral to a form containing the squares only of the variables. This method is applied in the memoir in question to the integrals of $n$ variables, analogous to those which give the attraction of an ellipsoid; and that directly without effecting an integration with respect to the radius vector. I proceed to show how the preceding investigations lead to Legendre's integral, and how the method in question effects with the utmost simplicity the integration which Legendre accomplished by means of what Poisson has spoken of as inextricable calculations.

Consider in particular the formula
\[
V = \int (x, y, \ldots)^{k}\, dx\, dy\ldots,
\]
the number of variables being as before $n$, and $(x, y, \ldots)^{k}$ denoting a homogeneous function of the order $k$. The equation for the limits is assumed to be
\[
l(x-a)^{2} + m(y-b)^{2} + \ldots = k.
\]
Assume
\[
\psi(x, y, \ldots) = \alpha(x^{2} + y^{2} + \ldots) + \beta(lx^{2} + my^{2} + \ldots) - (lax + mby + \ldots)^{2},
\]
(where $\delta = la^{2} + mb^{2} + \ldots - k$, is taken to be positive); or more simply,
\[
\psi(x, y, \ldots) = (\omega^{2} + l\delta - l^{2}a^{2})x^{2} + (\omega^{2} + m\delta - m^{2}b^{2})y^{2} + \ldots - 2lmab\,xy - \ldots
\]
Here $\mu = k + 2i - 3$. Also, putting for shortness
\[
lp + mq + \ldots = A, \qquad lp^{2} + mq^{2} + \ldots = \Phi,
\]
it is easy to obtain
\[
f\left(\frac{p}{\rho}, \frac{q}{\rho}, \ldots\right) = \frac{\rho^{k+2i+n-3}}{k + 2i + n - 3} \cdot \frac{(A + \omega\rho)^{k+2i+n-3} - (A - \omega\rho)^{k+2i+n-3}}{2A\rho^{k+2i+n-3}},
\]
values which give
\[
V = \frac{(-)^{i}}{k + 2i + n - 3} \int \frac{(A^{2} - \omega^{2}\rho^{2})^{i}\, \omega^{n-3}}{\Phi^{\frac{1}{2}(k+2i+n-3)}}\, [(A + \omega\rho)^{k+2i+n-3} - (A - \omega\rho)^{k+2i+n-3}]\, (p, q, \ldots)^{k}\, \mathcal{S}\, d\omega\, d\theta\ldots
\]

\newpage

\setcounter{page}{443}

where it will be remembered that $p$, $q$, \ldots\ are linear functions (with constant terms) of $(n-1)$ variables $\eta$, $\zeta$, \ldots, these last mentioned quantities being themselves functions of $(n-2)$ variables $\theta$, \&c.\ such that $1 - \eta^{2} - \zeta^{2} - \ldots = 0$ identically. If besides we suppose that $\Phi = lp^{2} + mq^{2} + \ldots$ reduces itself to the form $\mathfrak{a}\eta^{2} + \mathfrak{b}\zeta^{2} + \ldots$, we have, by the formula of the paper ``On the Simultaneous Transformation of two Homogeneous Equations of the Second Order,'' [74],
\[
(-\mathfrak{b})(-\mathfrak{c})\ldots =
\frac{(\omega^{2} + l\delta - l\lambda)}{(\omega^{2} + l\delta - l\lambda)} \cdot \frac{(\omega^{2} + m\delta - m\lambda)}{(\omega^{2} + m\delta - m\lambda)} \cdots,
\]
which is true, whatever be the value of $\lambda$.

It seems difficult to proceed further with the general formula, and I shall suppose $n = 3$, $i = 0$, $k = 1$, $(x, y, \ldots)^{k} = x$, or write
\[
V = \int \frac{x\, dx\, dy\, dz}{\sqrt{x^{2} + y^{2} + z^{2}}},
\]
the equation of the limits being
\[
l(x-a)^{2} + m(y-b)^{2} + n(z-c)^{2} = k.
\]
Here we may assume $\eta = \cos\theta$, $\zeta = \sin\theta$, (values which give $\mathcal{S} = 1$). And we have
\[
I = \frac{1}{2}\omega^{2} \int_{0}^{2\pi} \frac{(\alpha + \alpha'\cos\theta + \alpha''\sin\theta)\, d\theta}{\sqrt{\mathfrak{a}\cos^{2}\theta + \mathfrak{b}\sin^{2}\theta + 2\mathfrak{h}\cos\theta\sin\theta}},
\]
from $\theta = 0$ to $\theta = 2\pi$; or, what comes to the same thing,
\[
I = \omega^{2} \int_{0}^{\pi} \frac{(\alpha + \alpha'\cos\theta + \alpha''\sin\theta)\, d\theta}{\sqrt{\mathfrak{a}\cos^{2}\theta + \mathfrak{b}\sin^{2}\theta + 2\mathfrak{h}\cos\theta\sin\theta}}.
\]
Hence
\[
V = 4\pi \int \frac{\omega^{2}P\, d\omega}{\sqrt{(\mathfrak{B} - mP^{2})(\mathfrak{C} - nP^{2}) - \mathfrak{F}^{2}}},
\]
we have from the formul\ae\ of the paper before quoted,
\[
P = \frac{1}{\mathfrak{a}}\left[(\mathfrak{B} - mP^{2})(\mathfrak{C} - nP^{2}) - \mathfrak{F}^{2}\right] \cdot \left\{\frac{\mathfrak{B} - mP^{2}}{(\mathfrak{B} - mP^{2})(\mathfrak{C} - nP^{2}) - \mathfrak{F}^{2}}\right\}^{\frac{1}{2}} \cdot \ldots,
\]
$\mathfrak{B}$, $\mathfrak{C}$, $\mathfrak{F}$, being the coefficients of $y^{2}$, $z^{2}$, $yz$ in $\psi(x, y, z)$, viz.
\[
\mathfrak{B} = \omega^{2} + m\delta - m^{2}b^{2}, \quad \mathfrak{C} = \omega^{2} + n\delta - n^{2}c^{2}, \quad \mathfrak{F} = mnbc;
\]
and consequently
\[
V = 4\pi \int \frac{\omega^{2}P\, d\omega}{\sqrt{PQR}\, \sqrt{(\mathfrak{B} - mP^{2})(\mathfrak{C} - nP^{2}) - \mathfrak{F}^{2}}},
\]
where
\[
P = \frac{1}{\mathfrak{a}}\left[(\mathfrak{B} - mP^{2})(\mathfrak{C} - nP^{2}) - \mathfrak{F}^{2}\right].
\]

\newpage

\setcounter{page}{444}

Also from the equation
\[
\frac{l^{2}a^{2}}{\omega^{2} + l\delta - lP} + \frac{m^{2}b^{2}}{\omega^{2} + m\delta - mP} + \frac{n^{2}c^{2}}{\omega^{2} + n\delta - nP} = 1,
\]
differentiating with respect to $\lambda$, and writing $\lambda = P$,
\begin{multline*}
\frac{PQR}{K(P - Q)(P - R)} \\
= -\frac{1}{K}\left\{\frac{1}{(\omega^{2} + l\delta - lP)^{2}} + \frac{1}{(\omega^{2} + m\delta - mP)^{2}} + \frac{1}{(\omega^{2} + n\delta - nP)^{2}}\right\},
\end{multline*}
or, as this may be written,
\begin{multline*}
\frac{PQR}{K(P - Q)(P - R)} \\
= \frac{1}{(\omega^{2} + l\delta - lP)(\omega^{2} + m\delta - mP)(\omega^{2} + n\delta - nP)} \left\{\frac{1}{(\omega^{2} + l\delta - lP)^{2}} + \cdots\right\},
\end{multline*}
and from the values first written down, for $\mathfrak{B}$, $\mathfrak{C}$, $\mathfrak{F}$, we obtain $(\mathfrak{B} - mP^{2})(\mathfrak{C} - nP^{2}) - \mathfrak{F}^{2}$
\begin{align*}
&= (\omega^{2} + m\delta)(\omega^{2} + n\delta) - m^{2}b^{2}(\omega^{2} + n\delta) - n^{2}c^{2}(\omega^{2} + m\delta) \\
&\quad - mP(\omega^{2} + n\delta - n^{2}c^{2}) - nP(\omega^{2} + m\delta - m^{2}b^{2}) + mnP^{2} \\
&= (\omega^{2} + m\delta - mP)(\omega^{2} + n\delta - nP) - m^{2}b^{2}(\omega^{2} + n\delta - nP) - n^{2}c^{2}(\omega^{2} + m\delta - mP) \\
&= \frac{K}{P - Q}\cdot\frac{K}{P - R} \cdot P, \text{ the last reduction being effected by means of the equation}
\end{align*}
\[
\frac{l^{2}a^{2}}{\omega^{2} + l\delta - lP} + \frac{m^{2}b^{2}}{\omega^{2} + m\delta - mP} + \frac{n^{2}c^{2}}{\omega^{2} + n\delta - nP} = 1.
\]
Hence
\begin{multline*}
\frac{P^{4}\sqrt{QR}\, d\omega}{K(P - Q)(P - R)\sqrt{(\mathfrak{B} - mP^{2})(\mathfrak{C} - nP^{2}) - \mathfrak{F}^{2}}} \\
= \frac{la\, d\omega}{(\omega^{2} + l\delta - lP)^{\frac{1}{2}}(\omega^{2} + m\delta - mP)^{\frac{1}{2}}(\omega^{2} + n\delta - nP)^{\frac{1}{2}}} \cdot \frac{1}{\sqrt{\frac{l^{2}a^{2}}{(\omega^{2} + l\delta - lP)^{2}} + \frac{m^{2}b^{2}}{(\omega^{2} + m\delta - mP)^{2}} + \frac{n^{2}c^{2}}{(\omega^{2} + n\delta - nP)^{2}}}}.
\end{multline*}
Substituting this value, and multiplying out the fractions in the denominator,
\[
V = 4\pi la \int \frac{(\omega^{2} + l\delta - lP)^{\frac{1}{2}}(\omega^{2} + m\delta - mP)^{\frac{1}{2}}(\omega^{2} + n\delta - nP)^{\frac{1}{2}}\, d\omega}{l^{2}a^{2}(\omega^{2} + m\delta - mP)^{2}(\omega^{2} + n\delta - nP)^{2} + m^{2}b^{2}(\omega^{2} + n\delta - nP)^{2}(\omega^{2} + l\delta - lP)^{2} + n^{2}c^{2}(\omega^{2} + l\delta - lP)^{2}(\omega^{2} + m\delta - mP)^{2}},
\]
the reduction of which integral has been already treated of in the former part of this present memoir.



\newpage
% ============================================================
% PAPER 76
% ============================================================
\section*{76. On the triple tangent planes of surfaces of the third order.}
\markboth{On the triple tangent planes of surfaces of the third order.}{}
\addcontentsline{toc}{section}{76. On the triple tangent planes of surfaces of the third order.}

\begin{flushright}
\textit{[From the Cambridge and Dublin Mathematical Journal, vol.\ IV.\ (1849), pp.\ 118--132.]}
\end{flushright}

\setcounter{page}{445}

\textsc{A} \textsc{surface} of the third order contains in general a certain number of straight lines. Any plane through one of these lines intersects the surface in the line and in a conic, that is in a curve or system of the third order having two double points. Such a plane is therefore a double tangent plane of the surface, the double points (or points where the line and conic intersect) being the points of contact. By properly determining the plane, the conic will reduce itself to a pair of straight lines. Here the plane intersects the surface in three straight lines, that is in a curve or system of the third order having three double points, and the plane is therefore a triple tangent plane, the three double points or points of intersection of the lines taken two and two together being the points of contact. The number of lines and triple tangent planes is determined by means of a theorem very easily demonstrated, viz.\ that through each line there may be drawn five (and only five) triple tangent planes. Thus, considering any triple tangent plane, through each of the three lines in this plane there may be drawn (in addition to the plane in question) four triple tangent planes: these twelve new planes give rise to twenty-four new lines upon the surface, making up with the former three lines, twenty-seven lines upon the surface. It is clear that there can be no lines upon the surface besides these twenty-seven; for since the three lines upon the triple tangent plane are the complete intersection of this plane with the surface, every other line upon the surface must meet the triple tangent plane in a point upon one of the three lines, and must therefore lie in a plane passing through one of these lines, such plane (since it meets the surface in two lines and therefore in a third line) being obviously a triple tangent plane. Hence the whole number of lines upon the surface is twenty-seven; and it immediately follows that the number of triple tangent planes is forty-five. The number of lines upon the surface may also be obtained by the following method, which has the advantage of not assuming

\newpage

\setcounter{page}{446}

\textit{a priori} the existence of a line upon the surface. Imagine the cone having for its vertex a given point not upon the surface and circumscribed about the surface, every double tangent plane of the cone is also a double tangent plane of the surface, and therefore intersects the surface in a straight line (and a conic). And, conversely, if there be any line upon the surface, the plane through this line and the vertex of the cone will be a double tangent plane of the cone. Hence the number of double tangent planes of the cones is precisely that of the lines upon the surface. By the theorems in Mr [Dr] Salmon's paper ``On the degree of a surface reciprocal to a given one,'' Journal, vol.~II. [1847] p.~65, the cone is of the sixth order and has no double lines and six cuspidal lines: hence by the formula in Pl"{u}cker's ``Theorie der algebraischen Curven,'' [1839] p.~211, stated so as to apply to cones instead of plane curves, viz.\ $n$ being the order, $x$ the number of double lines, $y$ that of the cuspidal lines, $u$ that of the double tangent planes, then
\[
u = \tfrac{1}{2}n(n-2)(n^{2}-9) - (2n + 3y)(n^{2}-n-6) + 2x(x-1) + 6xy + \tfrac{1}{2}y(y-1),
\]
the number of double tangent planes is twenty-seven, which is therefore also the number of lines upon the surface.

Suppose the equation of one of the triple tangent planes to be $w = 0$, and let $x = 0$, $y = 0$, be the equation of any two triple tangent planes intersecting the plane $w = 0$ in two of the lines in which it meets the surface. Let $z = 0$ be the equation of a triple tangent plane meeting $w = 0$ in the remaining line in which it intersects the surface. The equation of the surface of the third order is in every case of the form $wP + kxyz = 0$, $P$ being a function of the second order, but of the four different planes which the equation $z = 0$ may be supposed to represent, one of them such that the function $P$ resolves itself into the product of a pair of factors, and for the remaining three this resolution into factors does not take place. This will be obvious from the sequel: at present I shall suppose that the plane $z = 0$ is of the latter class, or that $P = 0$ represents a proper surface of the second order. Since $x = 0$, $y = 0$, $z = 0$, are treble tangent planes of the surface, each of these planes must be a tangent plane of the surface of the second order $P = 0$, and this will be the case if we assume
\[
P = x^{2} + y^{2} + z^{2} + w^{2},
\]
and considering $x$, $y$, $z$ and $w$ as each of them implicitly containing an arbitrary constant, this is the most general function which satisfies the conditions in question.

We are thus led to the equation of the surface of the third order:

\newpage

\setcounter{page}{447}

\begin{center}
\textbf{Surfaces of the third order.}
\end{center}

I have found that by expressing the parameter $k$ in the particular form
\[
k = -\frac{1}{lmn}\left(\frac{p^{2} - \alpha^{2}}{2lmn}\right),
\]
or, as this equation may be more conveniently written,
\[
\alpha = \sqrt{lmn + \frac{1}{lmn}(p^{2} - \alpha^{2})},
\]
the equations of all the planes are expressible in a rational form. These equations are in fact the following: [I have added, here and in the table p.~450, the reference numbers $12'$, $23'$, \&c.\ constituting a different notation for the lines and planes.]

\begin{table}[h]
\centering
\begin{tabular}{@{}lll@{}}
\toprule
& Equation & Lines \\
\midrule
$(\alpha)$ & $w = 0$ & $12'$ \\
$(\beta)$ & $lx + my + nz + w = 0$ & $34'$ \\
$(\gamma)$ & $x = 0$ & $12.34.56$ \\
$(\delta)$ & $y = 0$ & $42'$ \\
$(\epsilon)$ & $z = 0$ & $14'.23'$ \\
$(\zeta)$ & $\frac{1}{m}x + \frac{1}{n}y + \frac{1}{l}z + w = 0$ & $31'$ \\
$(\eta)$ & $x + l(m - \alpha)y + \ldots = 0$ & $41'$ \\
$(\theta)$ & $y + m(\ldots) = 0$ & $34'$ \\
$(\iota)$ & $lx + \frac{1}{m}y + nz + w = 0$ & $13.24.56$ \\
$(\kappa)$ & $\frac{1}{l}x + my + \frac{1}{n}z + w = 0$ & $24'$ \\
$(\lambda)$ & $\frac{1}{l}x + \frac{1}{m}y + nz + w = 0$ & $23'$ \\
\bottomrule
\end{tabular}
\end{table}

A somewhat more elegant form is obtained by writing $p = 2q + \alpha$; this gives

\newpage

\setcounter{page}{448}

\begin{table}[h]
\centering
\begin{tabular}{@{}lll@{}}
\toprule
& Equation & Lines \\
\midrule
$(\iota)$ & $lx + \frac{1}{m}y + nz + w = 0$ & $24'$ \\
$(\kappa)$ & $\frac{1}{l}x + my + nz + w = 0$ & $14.23.56$ \\
$(\lambda)$ & $\frac{1}{l}x + my + \frac{1}{n}z + w = 0$ & $23'$ \\
$(\mu)$ & $\frac{m(p-\alpha) + 2nl}{p}x + y + \ldots = 0$ & $12.35.46$ \\
$(\nu)$ & $y + \frac{n(p-\alpha) + 2lm}{p}z + \ldots = 0$ & $52'$ \\
$(\xi)$ & $\frac{l(p-\alpha) + 2mn}{p}x + z + \ldots = 0$ & $12.36.45$ \\
$(o)$ & $x + \frac{m(p-\alpha) + 2nl}{p}y + \ldots = 0$ & $62'$ \\
$(\pi)$ & $z + \frac{n(p-\alpha) + 2lm}{p}x + \ldots = 0$ & $16'$ \\
$(\rho)$ & $y + \frac{n(p-\alpha)}{p}z + \ldots = 0$ & $56'$ \\
$(\sigma)$ & $\frac{1}{l}x + n(p-\alpha)y + nz + w = 0$ & $45'$ \\
$(\tau)$ & $mx + \frac{1}{n}y - \frac{m(p-\alpha)}{p}z + w = 0$ & $34'$ \\
$(\upsilon)$ & $-\frac{n(p-\alpha)}{p}x + my + \frac{1}{n}z + w = 0$ & $15.26.34$ \\
$(\phi)$ & $\frac{1}{m}x + y + nz + w = 0$ & $16.24.35$ \\
$(\chi)$ & $lx + \frac{n(p-\alpha)}{p}y - m(\frac{p-\alpha}{p}z) + w = 0$ & $14.25.36$ \\
$(\psi)$ & $-\frac{l(p-\alpha)}{p}x + y + mnz + w = 0$ & $65'$ \\
$(\omega)$ & $\frac{1}{m}x + \frac{l(p-\alpha)}{p}y + nz + w = 0$ & $46'$ \\
$(\alpha_{1})$ & $lx + \frac{1}{m}y - \frac{n(p-\alpha)}{p}z + w = 0$ & $4'6$ \\
\bottomrule
\end{tabular}
\end{table}

\newpage

\setcounter{page}{449}

\begin{table}[h]
\centering
\begin{tabular}{@{}lll@{}}
\toprule
& Equation & Lines \\
\midrule
$(\mu_{1})$ & $\frac{1}{l}x - \frac{m(p-\alpha)}{p}y + nz + w = 0$ & $16.25.34$ \\
$(\nu_{1})$ & $lx + \frac{1}{n}y - \frac{m(p-\alpha)}{p}z + w = 0$ & $15.24.36$ \\
$(\xi_{1})$ & $\frac{1}{l}x + \frac{1}{m}y - \frac{n(p-\alpha)}{p}z + w = 0$ & $14.26.35$ \\
$(p)$ & $\frac{1}{l}x + ny + mz + [mn(p-\alpha) - 2l(1-m^{2}-n^{2})]\frac{w}{p} = 0$ & $51'$ \\
$(q)$ & $nx + \frac{1}{m}y + lz + [nl(p-\alpha) - 2m(1-n^{2}-l^{2})]\frac{w}{p} = 0$ & $35'$ \\
$(r)$ & $mx + ly + \frac{1}{n}z + [lm(p-\alpha) - 2n(1-l^{2}-m^{2})]\frac{w}{p} = 0$ & $13.25.46$ \\
$(p')$ & $\frac{1}{l}x + \frac{1}{2}y + \frac{1}{2}z + \ldots = 0$ & $26'$ \\
$(q')$ & $\frac{1}{2}x + \frac{1}{m}y + \frac{1}{2}z + \ldots = 0$ & $16.23.45$ \\
$(r')$ & $\frac{1}{2}x + \frac{1}{2}y + \frac{1}{n}z + \ldots = 0$ & $63'$ \\
$(s')$ & $\frac{w}{p-\alpha} + x + \ldots = 0$ & $25'$ \\
$(t')$ & $\frac{w}{p-\alpha} + y + \ldots = 0$ & $15.23.46$ \\
$(u')$ & $\frac{w}{p-\alpha} + z + \ldots = 0$ & $53'$ \\
$(p_{1})$ & $\frac{w}{p-\alpha} + x + ny + mz + \ldots = 0$ & \\
$(q_{1})$ & $nx + \frac{w}{p-\alpha} + y + lz + \ldots = 0$ & $36'$ \\
$(r_{1})$ & $mx + ly + \frac{w}{p-\alpha} + z + \ldots = 0$ & $13.26.45$ \\
\bottomrule
\end{tabular}
\end{table}

In fact, representing the several functions on the left-hand side of these equations respectively by the letters placed opposite to them respectively, the function $U$ is expressible in the sixteen forms following
\begin{align*}
U &= w\alpha\alpha_{1} + k\xi\eta\zeta, \\
&= w\beta\beta_{1} + k\eta\zeta\xi, \\
&= w\gamma\gamma_{1} + k\zeta\xi\eta, \\
&= w\delta\delta_{1} + k\xi\eta\zeta,
\end{align*}
(being the forms containing $w$, out of a complete system of one hundred and twenty different forms).

\newpage

\setcounter{page}{450}

\begin{align*}
U &= w\epsilon\epsilon_{1} + k\eta\zeta x, \\
&= w\zeta\zeta_{1} + k\zeta\xi y, \\
&= w\eta\eta_{1} + k\xi\eta z, \\
&= w\iota\iota_{1} + k\eta\zeta x, \\
&= w\kappa\kappa_{1} + k\zeta\xi y, \\
&= w\lambda\lambda_{1} + k\xi\eta z, \\
&= w\mu\mu_{1} + k\eta\zeta x, \\
&= w\nu\nu_{1} + k\zeta\xi y, \\
&= w\xi\xi_{1} + k\xi\eta z, \\
&= w\pi\pi_{1} + k\eta\zeta x, \\
&= w\rho\rho_{1} + k\zeta\xi y, \\
&= w\sigma\sigma_{1} + k\xi\eta z.
\end{align*}

The forty-five planes pass five and five through the twenty-seven lines in the following manner:
\begin{align*}
&(\alpha_{1})\ (w, x, \xi, \iota, \varkappa) \ldots 12, \quad &&(\alpha_{4})\ (x, \varrho, h, \iota, \iota_{1}) \ldots 34, \quad &&(\alpha_{7})\ (x, m_{1}, n, q_{1}, r) \ldots 56, \\
&(\beta_{1})\ (w, y, \eta, \nu, \lambda) \ldots 16, \quad &&(\beta_{4})\ (y, h, f, \nu, \nu_{1}) \ldots 24, \quad &&(\beta_{7})\ (y, n_{1}, \iota, r_{1}, p) \ldots 6, \\
&(\gamma_{1})\ (w, z, \zeta, \mu, \lambda) \ldots 1, \quad &&(\gamma_{4})\ (z, f, g, n, n_{1}) \ldots 14, \quad &&(\gamma_{7})\ (z, \iota, m, p_{1}, q) \ldots 45, \\
&(\delta_{1})\ (f, f, e, p, p_{1}) \ldots 2, \quad &&(\delta_{4})\ (x, g, h, \iota, \iota_{1}) \ldots 56, \quad &&(\delta_{7})\ (x, m_{1}, n, q_{1}, r) \ldots 12, \\
&(\epsilon_{1})\ (\eta, g, \xi, q, q_{1}) \ldots 23, \quad &&(\epsilon_{4})\ (y, h, f, m, m_{1}) \ldots 4, \quad &&(\epsilon_{7})\ (y, n_{1}, \iota, r_{1}, p) \ldots 25, \\
&(\zeta_{1})\ (\zeta, K, \theta, r, r_{1}) \ldots 3, \quad &&(\zeta_{4})\ (z, f, g, n, n_{1}) \ldots 4', \quad &&(\zeta_{7})\ (z, \iota, m, p_{1}, q) \ldots 5',
\end{align*}
where each line may be represented by the letter placed opposite to the system of planes passing through it. The twenty-seven lines lie three and three upon the forty-five planes in the following manner:
\begin{align*}
&(w)\ \alpha_{1}\beta_{1}\gamma_{1}, \quad &&(f)\ \alpha_{4}\beta_{4}\gamma_{4}, \quad &&(\iota)\ \alpha_{7}\beta_{7}\gamma_{7}, \quad &&(p)\ \delta_{1}\epsilon_{1}\zeta_{1}, \\
&(x)\ \alpha_{4}\beta_{4}\gamma_{4}, \quad &&(g)\ \delta_{4}\epsilon_{4}\zeta_{4}, \quad &&(m)\ \delta_{7}\epsilon_{7}\zeta_{7}, \quad &&(q)\ \delta_{4}\epsilon_{4}\zeta_{4}, \\
&(y)\ \alpha_{1}\beta_{1}\gamma_{1}, \quad &&(h)\ \epsilon_{4}\zeta_{4}\alpha_{4}, \quad &&(n)\ \zeta_{4}\alpha_{4}\beta_{4}, \quad &&(r)\ \zeta_{4}\alpha_{4}\beta_{4}.
\end{align*}

\newpage

\setcounter{page}{451}

\begin{align*}
&(x)\ \alpha_{1}\alpha_{4}\alpha_{7}, \quad &&(\xi)\ \alpha_{4}\beta_{4}\gamma_{4}, \quad &&(\iota)\ \alpha_{7}\beta_{7}\gamma_{7}, \\
&(p)\ \alpha_{1}\beta_{4}\gamma_{7}, \quad &&(m)\ \beta_{4}\gamma_{7}\alpha_{1}, \\
&(y)\ \beta_{1}\beta_{4}\beta_{7}, \quad &&(g)\ \beta_{4}\gamma_{4}\alpha_{4}, \quad &&(q)\ \beta_{7}\gamma_{7}\alpha_{7}, \\
&(h)\ \gamma_{4}\alpha_{4}\beta_{4}, \quad &&(n)\ \alpha_{4}\beta_{4}\gamma_{4}, \quad &&(r)\ \gamma_{7}\alpha_{7}\beta_{7}, \\
&(f)\ \alpha_{1}\beta_{4}\alpha_{7}, \quad &&(x)\ \alpha_{7}\beta_{7}\gamma_{7}, \\
&(l)\ \alpha_{4}\beta_{7}\gamma_{7}, \quad &&(p_{1})\ \alpha_{1}\beta_{4}\gamma_{7}, \\
&(y)\ \beta_{1}\beta_{4}\beta_{7}, \quad &&(m_{1})\ \beta_{4}\gamma_{7}\alpha_{7}, \\
&(z)\ \gamma_{1}\gamma_{4}\gamma_{7}, \quad &&(n_{1})\ \gamma_{4}\alpha_{4}\beta_{4}, \quad &&(r_{1})\ \gamma_{7}\alpha_{7}\beta_{7}.
\end{align*}

The preceding method was the one that first occurred to me, and which appears to conduct most simply to the actual analytical expressions for the forty-five planes; but it is worth noticing that the relations between the lines and planes might have been obtained almost without algebraical developments, if we had supposed that $P$, instead of representing a proper surface of the second order, had represented a pair of planes. This would have conducted at once to one of the one hundred and twenty forms $U$, e.g.\ $U = wdd_{1} + k\xi\eta\zeta$. Or changing the notation so as to include $k$ in one of the linear functions, $U = ace - bdf$, and it is indeed obvious \textit{a priori}, by merely reckoning the number of arbitrary constants, that any function of the third order can be put under this form. If we suppose $a = \mu b$ to be the equation of one of the triple tangent planes through the intersection of the planes $a$ and $b$, the plane $a = \mu b$ meets the surface in the same lines in which it meets the hyperboloid $ace - df = 0$, that is, the two lines in the plane are generating lines of different species, and consequently one of them meets the pair of lines $cd$ and $ef$, and the other of them meets the pair of lines $cf$ and $de$ (where $cd$ represents the line of intersection of the planes $c = 0$, $d = 0$, \&c.). This suggests a notation for the lines in question, viz.\ each line may be represented by the three lines which it meets, or by the symbols $ab.cd.ef$ and $ab.cf.de$. Or observing that $\mu$ has three values, and that the same considerations apply \textit{mutatis mutandis} to the planes through $bc$ and $ca$, the whole system of lines may be represented by the notation
\begin{align*}
&(ab.cd.ef)_{1}, \quad (ab.cd.ef)_{2}, \quad (ab.cd.ef)_{3}, \\
&(ad.cf.eb)_{1}, \quad (ad.cf.eb)_{2}, \quad (ad.cf.eb)_{3}, \\
&(af.cb.ed)_{1}, \quad (af.cb.ed)_{2}, \quad (af.cb.ed)_{3}, \\
&(ab.cf.ed)_{1}, \quad (ab.cf.ed)_{2}, \quad (ab.cf.ed)_{3}, \\
&(ad.cb.ef)_{1}, \quad (ad.cb.ef)_{2}, \quad (ad.cb.ef)_{3}, \\
&(af.cd.eb)_{1}, \quad (af.cd.eb)_{2}, \quad (af.cd.eb)_{3},
\end{align*}
where the last eighteen lines have been divided into two systems of nine each.

\newpage

\setcounter{page}{452}

The five planes through $(ab.cd.ef)_{1}$ may be considered as cutting the surface in
\begin{align*}
&ab; \quad (ab.cf.ed)_{1}, \\
&cd; \quad (af.cd.eb)_{1}, \\
&ef; \quad (ad.cb.ef)_{1}, \\
&(ad.cf.eb)_{2}; \quad (af.cb.ed)_{2}, \\
&(ad.cf.eb)_{3}; \quad (af.cb.ed)_{3},
\end{align*}
(which supposes however that the distinguishing suffixes $1$, $2$, $3$, are added to the different planes according to a certain rule). And similarly for the lines in the planes through the other lines represented by symbols of the like form. The five planes through $ab$ intersect the surface in the lines
\begin{align*}
&(ab.cd.ef)_{1}, \quad (ab.cf.ed)_{1}, \\
&(ab.cd.ef)_{2}, \quad (ab.cf.ed)_{2}, \\
&(ab.cd.ef)_{3}, \quad (ab.cf.ed)_{3},
\end{align*}
and similarly for the planes through the other lines represented by symbols of a like form.

Observing that $\xi$, $\eta$, $\zeta$, correspond to $b$, $d$, $f$, and $w$, $\theta$, $\Theta$, to $a$, $c$, $e$ respectively, $ab$ corresponds to the intersection of $w$ and $\xi$, i.e.\ to $a_{1}$, \&c.; also $(ab.cd.ef)_{1}$, $(ab.cd.ef)_{2}$, $(ab.cd.ef)_{3}$ correspond to three lines meeting $a_{1}$, $b_{1}$, and $c_{1}$, i.e.\ to $\alpha_{1}$, $\alpha_{2}$, $\alpha_{3}$, \&c.; and the system of the twenty-seven lines as last written down corresponds to the system,
\begin{align*}
&a_{1}, a_{2}, a_{3}, \\
&a_{5}, a_{7}, a_{9}, \\
&b_{1}, b_{7}, b_{8}, \\
&c_{5}, c_{7}, c_{9}, \\
&a_{4}, a_{6}, a_{8}, \\
&b_{4}, b_{6}, b_{3}, \\
&c_{4}, c_{6}, c_{8}.
\end{align*}

The investigations last given are almost complete in themselves as the geometrical theory of the subject: there is however some difficulty in seeing \textit{a priori} the nature of the correspondence between the planes which determines which are the planes which ought to be distinguished with the same one of the symbolic numbers, $1$, $2$, $3$.

\newpage

\setcounter{page}{453}

There is great difficulty in conceiving the complete figure formed by the twenty-seven lines, indeed this can hardly I think be accomplished until a more perfect notation is discovered. In the mean time it is easy to find theorems which partially exhibit the properties of the system. For instance, any two lines, $a_{1}$, $b_{2}$, which do not meet are intersected by five other lines, $a_{3}$, $b_{4}$, $K_{3}$, $O_{3}$, $a_{5}$, (no two of which meet). Any four of these last-mentioned lines are intersected by the lines $a_{1}$, $b_{2}$ and no other lines, but any three of them, e.g.\ $a_{3}$, $a_{4}$, $a_{5}$, are intersected by the lines $a_{1}$, $b_{2}$, and by some third line (in the case in question the line $c_{3}$). Or generally any three lines, no two of which meet, are intersected by three other lines, no two of which meet.

Again, the lines which do not meet any one of the lines $a_{1}$, $a_{2}$, $a_{3}$, are $a_{4}$, $a_{5}$, $b_{3}$, $b_{4}$, $c_{1}$, $c_{2}$: these lines form a hexagon, the pairs of the opposite sides of which, $a_{4}$, $b_{3}$; $a_{5}$, $c_{1}$; $b_{4}$, $c_{2}$, are met by the pairs $a_{1}$, $b_{2}$; $c_{3}$, $a_{2}$ and $a_{3}$, $b_{2}$, respectively, viz.\ by pairs out of the system of three lines intersecting the system $a_{1}$, $a_{2}$, $a_{3}$. And the lines $a_{1}$, $a_{2}$, $a_{3}$ may be considered as representing any three lines no two of which meet. Again, consider three lines in the same triple tangent plane, e.g.\ $a_{1}$, $b_{1}$, $c_{1}$, and the hexahedron formed by any six triple tangent planes passing two and two through these lines, e.g.\ the planes $x$, $y$, $z$, $\xi$, $\eta$, $\zeta$. These planes contain (independently of the lines $a_{1}$, $b_{1}$, $c_{1}$) the twelve lines $a_{2}$, $a_{3}$, $a_{4}$, $a_{5}$, $b_{2}$, $b_{3}$, $b_{4}$, $b_{5}$, $c_{2}$, $c_{3}$, $c_{4}$, $c_{5}$. Consider three contiguous faces of the hexahedron, e.g.\ $x$, $y$, $z$, the lines in these planes, viz.\ $a_{2}$, $b_{5}$, $d$, $a_{3}$, $b_{4}$, $c_{2}$, form a hexagon the opposite sides of which intersect in a point, or in other words these six lines are generating lines of a hyperboloid. The same property holds for the systems $\xi$, $\eta$, $\zeta$; $\xi$, $y$, $\zeta$; $x$, $\eta$, $\zeta$; $x$, $y$, $\xi$; but for the system $\xi$, $\eta$, $\zeta$, the six lines are $a_{2}$, $b_{3}$, $c_{4}$, and $a_{3}$, $b_{4}$, $c_{5}$, which form two triangles, and similarly for the systems $\xi$, $y$, $z$; $x$, $\eta$, $z$; and $x$, $y$, $\xi$; so that the twelve lines form four hexagons (the opposite sides of which intersect) circumscribed round four of the angles of the hexahedron, and four pairs of triangles about the opposite four angles of the hexahedron. The number of such theorems might be multiplied indefinitely, and the number of different combinations of lines or planes to which each theorem applies is also very considerable.

Consider the four planes $x$, $\xi$, $\chi$, $\varkappa$, and represent for a moment the equations of these planes by $x + Aw = 0$, $x + Bw = 0$, $x + Cw = 0$, $x + Dw = 0$, so that
\[
A = 0, \quad B = \frac{1}{l}\sqrt{\frac{m}{l}\frac{n}{l}}, \quad C = l - (p - \alpha)\frac{1}{l} + 2mn, \quad D = \ldots
\]
By the assistance of
\[
B - C = \frac{[ln(p - \alpha) + 2m][lm(p - \alpha) + 2n]}{lmn(p - \alpha)}
\]
it is easy to obtain
\[
\frac{(A - C)(D - B)}{(A - D)(B - C)} = \frac{p - \beta}{p + \beta} \cdot \frac{[l(p - \alpha) + 2mn][m(p - \alpha) + 2nl][n(p - \alpha) + 2lm]}{[mn(p - \alpha) + 2l][nl(p - \alpha) + 2m][lm(p - \alpha) + 2n]},
\]

\newpage

\setcounter{page}{454}

which remains unaltered for cyclical permutations of $l$, $m$, $n$, i.e.\ the anharmonic ratio of $x$, $\xi$, $\chi$, $\varkappa$ is the same as that of $y$, $\eta$, $y'$, $y''$, or $z$, $\zeta$, $z'$, $z''$; there is of course no correspondence of $x$ to $y$ or $\xi$ to $\eta$, \&c., the correspondence is by the general properties of anharmonic ratios, a correspondence of the system $x$, $\xi$, $\chi$, $\varkappa$, to any one of the systems $(y, \eta, y', y'')$, or $(\eta, y, y', y'')$, or $(y, y', y'', \eta)$, or $(y, y', \eta, y'')$, indifferently.

The theorems may be stated generally as follows: ``Considering two lines in the same triple tangent plane, the remaining triple tangent planes through these two lines respectively are homologous systems.''

Suppose the surface of the third order intersected by an arbitrary plane. The curve of intersection is of course one of the third order, and the positions upon this curve of six of the points in which it is intersected may be arbitrarily assumed. Let these points be the points in which the plane is intersected by the lines $a_{1}$, $b_{1}$, $a_{2}$, $b_{2}$, $c_{1}$, $a_{3}$; or as we may term them, the points $a_{1}$, $b_{1}$, $a_{2}$, $b_{2}$, $c_{1}$, $a_{3}$. The point $c_{2}$ is of course the point in which the line $a_{1}b_{1}$ intersects the curve. The straight lines $a_{2}b_{2}$, $b_{3}c_{3}$, $c_{1}a_{1}$, and $a_{3}b_{3}$, $b_{1}c_{1}$, $c_{2}a_{2}$, show that $d$ and $b_{4}$ are the points in which $a_{2}b_{2}$ and $c_{1}c_{3}$ intersect the curve, and then $b_{3}$ and $c_{3}$ are determined as the intersections of $a_{3}c_{1}$, $a_{2}b_{1}$ with the curve. The intersection of the lines $b_{4}c_{4}$ and $a_{3}c_{3}$ (which is known to be a point upon the curve by the theorem, every curve of the third order passing through eight of the points of intersection of two curves of the third order passes through the ninth point of intersection) is the point $a_{4}$. The systems $a_{1}$, $b_{4}$, $c_{4}$; $a_{2}$, $b_{2}$, $c_{1}$; $a_{3}$, $b_{3}$, $c_{3}$, determine the conjugate system $a_{4}$, $b_{5}$, $c_{2}$; $a_{5}$, $b_{4}$, $c_{3}$, $a_{6}$, $b_{6}$, $c_{4}$; by reason of the straight lines $a_{1}a_{4}a_{5}$, $b_{1}b_{4}b_{5}$, $c_{1}c_{4}c_{5}$; $a_{1}a_{2}a_{3}$, $b_{1}b_{2}b_{3}$, $c_{1}c_{2}c_{3}$; $a_{1}a_{6}a_{7}$, $b_{1}b_{6}b_{7}$, $c_{1}c_{6}c_{7}$, viz.\ $a_{6}$ is the point where $a_{1}a_{4}$ intersects the curve, \&c. The relations of the systems $(a_{1}, b_{4}, c_{4}; a_{5}, b_{5}, c_{2})$, $(a_{3}, b_{3}, c_{3}; a_{7}, b_{7}, c_{7})$, $(a_{6}, b_{6}, c_{6}; a_{7}, b_{7}, c_{7})$ to the system $a_{1}$, $b_{1}$, $c_{1}$; $a_{2}$, $b_{2}$, $c_{2}$; $a_{3}$, $b_{3}$, $c_{3}$ are precisely identical. It is only necessary to show how the points $a_{5}$, $b_{5}$, $c_{5}$; $a_{6}$, $b_{6}$, $c_{6}$ of the latter system are determined by means of one of the former systems, suppose the system $a_{3}$, $b_{3}$, $c_{3}$; $a_{7}$, $b_{7}$, $c_{7}$; and to discover a compendious statement of the relation between the two systems. The points $a_{1}$, $b_{4}$, $c_{4}$; $a_{5}$, $b_{5}$, $c_{5}$; $a_{6}$, $b_{3}$, $c_{3}$; $a_{6}$, $b_{6}$, $c_{6}$; $a_{7}$, $b_{7}$, $c_{7}$, are a system of fifteen points lying on the fifteen straight lines $a_{1}b_{1}c_{1}$, $a_{2}b_{2}c_{2}$, $a_{3}b_{3}c_{3}$, $a_{1}a_{2}a_{3}$, $b_{1}b_{2}b_{3}$, $c_{1}c_{2}c_{3}$, $a_{1}a_{4}a_{5}$, $b_{1}b_{4}b_{5}$, $c_{1}c_{4}c_{5}$, $a_{2}b_{7}c_{4}$, $b_{2}a_{5}c_{6}$, $c_{2}a_{3}b_{5}$, $a_{3}b_{4}c_{4}$, $b_{3}c_{5}a_{6}$, $c_{3}a_{4}b_{6}$, viz.\ the nine points $a_{1}$, $b_{1}$, $c_{1}$; $a_{2}$, $b_{2}$, $c_{2}$; $a_{3}$, $b_{3}$, $c_{3}$ are the points of intersection of the three lines $a_{1}b_{1}c_{1}$, $a_{2}b_{2}c_{2}$, $a_{3}b_{3}c_{3}$ with the three lines $a_{1}a_{2}a_{3}$, $b_{1}b_{2}b_{3}$, $c_{1}c_{2}c_{3}$, and the remaining six points form a hexagon $a_{4}b_{7}c_{4}b_{5}a_{6}c_{6}$ of which the diagonals $a_{4}a_{5}$, $b_{7}b_{5}$, $c_{4}c_{6}$ pass through the points $a_{1}$, $b_{1}$, $c_{1}$, respectively, the alternate sides $a_{4}b_{7}$, $c_{6}a_{5}$, and $b_{4}c_{7}$ pass through the points $c_{2}$, $b_{1}$, $a_{3}$ respectively, and the remaining alternate sides $b_{4}c_{6}$, $c_{4}b_{5}$, and $b_{7}a_{6}$ pass through the three points $a_{1}$, $b_{2}$, $c_{3}$ respectively. The fifteen points of such a system do not necessarily lie upon a curve of the third order, as will presently be seen: in the actual case however where all the points lie upon a given curve of the third order, and the points $a_{1}$, $b_{1}$, $c_{1}$; $a_{2}$, $b_{2}$, $c_{2}$; $a_{7}$, $b_{7}$, $c_{7}$ are known, $a_{5}$, $b_{2}$, $c_{2}$; $a_{6}$, $b_{3}$, $c_{3}$ are the intersections of the curve with $a_{7}c_{1}$, $c_{2}a_{5}$, $a_{3}b_{4}$, $b_{7}a_{5}$, $c_{1}c_{3}$, $c_{4}b_{5}$ respectively, and the fact of the existence of the lines $a_{1}b_{4}c_{4}$, $a_{2}b_{4}c_{4}$, $a_{3}b_{4}c_{4}$, $a_{1}a_{2}a_{3}$, $b_{1}b_{2}b_{3}$, $c_{1}c_{2}c_{3}$ is an immediate consequence of the theorem quoted above with respect to curves of the third order---a theorem from which the entire system of relations between the twenty-seven points on the curve might have been deduced \textit{a priori}. But returning to the system of fifteen points, suppose the lines $a_{1}b_{1}c_{1}$, $a_{2}b_{2}c_{2}$, $a_{3}b_{3}c_{3}$, and $a_{1}a_{2}a_{3}$, $b_{1}b_{2}b_{3}$, and also the point $a_{4}$ to be given arbitrarily. The point $a_{5}$ lies on the line $a_{4}a_{6}$, suppose its position upon this line to be arbitrarily assumed (in which case, since the ten points $a_{1}$, $a_{2}$, $a_{3}$, $b_{1}$, $b_{2}$, $b_{3}$, $c_{1}$, $c_{2}$, $c_{3}$, are sufficient to determine a curve of the third order, there is no curve of the third order through these points and the point $a_{4}$). If the points $b_{7}$, $c_{7}$, $b_{4}$, $c_{4}$ can be so determined that the sides of the quadrilateral $b_{7}b_{4}c_{6}c_{4}$, viz.\ $b_{7}b_{4}$, $b_{7}c_{6}$, $c_{6}c_{4}$, $c_{4}b_{4}$ pass through the points $b_{1}$, $a_{3}$, $c_{1}$, $a_{2}$, respectively, while the angles $b_{7}$, $b_{4}$, $c_{6}$, $c_{4}$ lie upon the lines $a_{3}c_{1}$, $a_{3}c_{2}$, $a_{7}b_{4}$, and $a_{4}b_{3}$ respectively, the required conditions will be satisfied by the fifteen points in question; and the solution of this problem is known. I have not ascertained whether in the case of an arbitrary position as above of the point $a_{5}$, it is possible to determine a complete system of twenty-seven points lying three and three upon forty-five lines in the same manner as the twenty-seven points upon the curve of the third order; but it appears probable that this is the case, and to determine whether it be so or not, presents itself as an interesting problem for investigation.

\newpage

\setcounter{page}{455}

Suppose that the intersecting plane coincides with one of the triple tangent planes. Here we have a system of twenty-four points, lying eight and eight in three lines; the twenty-four points lie also three and three in thirty-two lines, which last-mentioned lines therefore pass four and four through the twenty-four points. If we represent by $a$, $b$, $c$, $d$, $a'$, $b'$, $c'$, $d'$ and $a$, $b$, $c$, $d$, $a'$, $b'$, $c'$, $d'$, the eight points, and eight points

\newpage

\setcounter{page}{456}

which lie upon two of the three lines (the order being determinate), the systems of four lines which intersect in the eight points of the third line are
\begin{align*}
&(aa', bb', cc', dd'), \quad &&(a'b', c'c', d'd'), \\
&(ab', ba', c'd, d'c), \quad &&(a'b', b'a, cd', c'd), \\
&(ac', ca', db', b'd), \quad &&(a'c, c'a, bd', d'b'), \\
&(ad', da', b'c, c'b), \quad &&(a'd, d'a, bc', cb'),
\end{align*}
the principle of symmetry made use of in this notation (which however represents the actual symmetry of the system very imperfectly) being obviously entirely different from that of the case of an arbitrary intersecting plane. The transition case where the intersecting plane passes through one of the lines upon the surface (and is thus a double tangent plane) would be worth examining. It should be remarked that the preceding theory is very materially modified when the surface of the third order has one or more conical points; and in the case of a double line (for which the surface becomes a ruled surface) the theory entirely ceases to be applicable. I may mention in conclusion that the whole subject of this memoir was developed in a correspondence with Mr Salmon, and in particular, that I am indebted to him for the determination of the number of lines upon the surface and for the investigations connected with the representation of the twenty-seven lines by means of the letters $a$, $c$, $e$, $b$, $d$, $f$, as developed above.



\newpage
% ============================================================
% PAPER 77
% ============================================================
\section*{77. On the order of certain systems of algebraical equations.}
\markboth{On the order of certain systems of algebraical equations.}{}
\addcontentsline{toc}{section}{77. On the order of certain systems of algebraical equations.}

\begin{flushright}
\textit{[From the Cambridge and Dublin Mathematical Journal, vol.\ IV.\ (1849), pp.\ 132--137.]}
\end{flushright}

\setcounter{page}{457}

\textsc{Suppose} the variables $x$, $y$, \ldots\ so connected that any one of the ratios $x : y : z$, \ldots\ or, more generally, any determinate function of these ratios, depends on an equation of the $\mu^{\text{th}}$ order. The variables $x$, $y$, $z$, \ldots\ are said to form a system of the $\mu^{\text{th}}$ order.

In the case of two variables $x$, $y$, supposing that these are connected by an equation $U = 0$ ($U$ being a homogeneous function of the order $\mu$) the variables form a system of the $\mu^{\text{th}}$ order; and, conversely, whenever the variables form a system of the $\mu^{\text{th}}$ order, they are connected by an equation of the above form.

In the case of a greater number of variables, the question is one of much greater difficulty. Thus with three variables $x$, $y$, $z$; if $\mu$ be resolvable into the factors $\mu'$, $\mu''$, then, supposing the variables to be connected by the equations $U = 0$, $V = 0$, $U$ and $V$ being homogeneous functions of the orders $\mu'$, $\mu''$, respectively, they will it is true form a system of the $\mu^{\text{th}}$ order, but the converse proposition does not hold: for instance, if $\mu$ is a prime number, the only mode of forming a system of the $\mu^{\text{th}}$ order would on the above principle be to assume $\mu = \mu'$, $\mu'' = 1$, that is to suppose the variables connected by an equation of the $\mu^{\text{th}}$ order and a linear equation; but this is far from being the most general method of obtaining such a system. In fact, systems not belonging to the class in question may be obtained by the introduction of subsidiary variables to be eliminated: the simplest example is the following: suppose $a$, $b$, $a'$, $b'$, $a''$, $b''$ to be linear functions (without constant terms) of $x$, $y$, $z$, and write
\begin{align*}
a\xi + b\eta &= 0, \\
a'\xi + b'\eta &= 0, \\
a''\xi + b''\eta &= 0;
\end{align*}
equations from which, by the elimination of $\xi$, $\eta$, two relations may be obtained between the variables $x$, $y$, $z$.

Suppose, however, from these three equations $x$, $y$, $z$ are first eliminated: the ratio $\xi : \eta$ will evidently be determined by a cubic equation; and assuming $\xi : \eta$ to be equal to one of the roots of this, any two of the three equations may be considered as implying the third; and will likewise determine linearly the ratios $x : y : z$. Hence any determinate function of these ratios depends on a cubic equation only, or the system is one of the third order. But the order of the system may be obtained by means of the equations resulting from the elimination of $\xi$, $\eta$; and since this will explain the following more general example (in which the corresponding process is the only one which readily offers itself), it will be convenient to deduce the preceding result in this manner. Thus, performing the elimination, we have
\[
L = (a'b'' - a''b') = 0, \quad L' = (a''b - ab'') = 0, \quad L'' = (ab' - a'b) = 0.
\]
Here the equations $L = 0$, $L' = 0$, $L'' = 0$, are each of them of the second order, and any two of them may be considered as implying the third. For we have identically,
\[
aL + a'L' + a''L'' = 0,\footnote{Also $bL + b'L' + b''L'' = 0$: but since by the elimination of $L''$, taking into account the actual values of $L$ and $L'$, we obtain an identical equation, these two relations may be considered as equivalent to a single one.}
\]
so that $L = 0$, $L' = 0$, gives $a''L'' = 0$, or $L'' = 0$. Nevertheless the system is imperfectly represented by means of two equations only. For instance, $L = 0$, $L' = 0$ do, of themselves, represent a system which is really of the fourth order. In fact, these equations are satisfied by $a'' = 0$, $b'' = 0$, (which is to be considered as forming a system of the first order), but these values do not satisfy the remaining equation $L'' = 0$. In other words, the equations $L = 0$, $L' = 0$ contain an extraneous system of the first order, and which is seen to be extraneous by means of the last equation $L''$: the system required is the system of the third order which is common to the three equations $L = 0$, $L' = 0$, $L'' = 0$.

\newpage

\setcounter{page}{458}

Suppose, more generally, that $x$, $y$, $z$ are connected by $p + 1$ equations, involving $p$ variables $\xi$, $\eta$, $\zeta$, \ldots,
\begin{align*}
a\xi + b\eta + c\zeta + \ldots &= 0, \\
a'\xi + b'\eta + c'\zeta + \ldots &= 0, \\
&\ \,\vdots
\end{align*}
or what comes to the same by the equations (equivalent to two independent relations)
\[
\begin{vmatrix}
a, & a', & a'', & \ldots & a^{(p)} \\
b, & b', & b'', & \ldots & b^{(p)} \\
\vdots & \vdots & \vdots & \ddots & \vdots
\end{vmatrix} = 0,
\]
(where the number of horizontal rows is $p$). Consider $x$, $y$, $z$, as connected by the two equations $= 0$:
\[
\begin{vmatrix}
a, & a', & \ldots & a^{(p-1)}, & a^{(p)} \\
b, & b', & \ldots & b^{(p-1)}, & b^{(p)} \\
\vdots & \vdots & \ddots & \vdots & \vdots
\end{vmatrix} = 0,
\]
these form a system of the order $p^{2}$, but they involve the extraneous system

\newpage

\setcounter{page}{459}

Suppose $\phi(p)$ is the order of the system in question, then the order of this last system is $\phi(p-1)$ and hence $\phi(p) = p^{2} - \phi(p-1)$: observing that $\phi(2) = 3$, this gives directly $\phi(p) = \frac{1}{2}p(p+1)$. Hence the order of the system is $\frac{1}{2}p(p+1)$.

Suppose $x$, $y$, $z$, connected by equations of the form $U = 0$, $V = 0$, $W = 0$; $U$, $V$, $W$ being linear in $x$, $y$, $z$, and homogeneous functions of the orders $m$, $n$, $p$ respectively in $\xi$, $\eta$. By eliminating $x$, $y$, $z$, the ratio $\xi : \eta$ will be determined by an equation of the order $m + n + p$; and since when this is known the ratios $x : y : z$ are linearly determinable, we have $m + n + p$ for the order $\mu$ of the system.

Thus, if $m = n = p = 2$, selecting the particular system
\begin{align*}
a\xi^{2} + 2b\xi\eta + c\eta^{2} &= 0, \\
b\xi^{2} + 2c\xi\eta + d\eta^{2} &= 0, \\
c\xi^{2} + 2d\xi\eta + e\eta^{2} &= 0,
\end{align*}
it is possible in this case to obtain two resulting equations of the orders two and three respectively, and which consequently constitute the system of the sixth order without containing any extraneous system. In fact, from the identical equation
\begin{multline*}
e^{2} \cdot (a\xi^{2} + 2b\xi\eta + c\eta^{2}) - (d\xi^{2} + 2e\xi\eta)(b\xi^{2} + 2c\xi\eta + d\eta^{2}) \\
+ (c\xi^{2} + 2d\xi\eta)(c\xi^{2} + 2d\xi\eta + e\eta^{2}) = (ae - 4bd + 3c^{2})\xi^{4},
\end{multline*}
and
\begin{multline*}
+(cd - be)(b\xi^{2} + 2c\xi\eta + d\eta^{2}) + (bd - c^{2})(c\xi^{2} + 2d\xi\eta + e\eta^{2}) \\
= (ace - ad^{2} - b^{2}e - c^{3} + 2bcd)\xi^{3}\eta,
\end{multline*}
we deduce
\begin{align*}
ae - 4bd + 3c^{2} &= 0, \\
ace - ad^{2} - b^{2}e - c^{3} + 2bcd &= 0;
\end{align*}
which form the system in question, and may for shortness be represented by $I = 0$, $J = 0$.

The three equations in $\xi$, $\eta$ may be considered as expressing that
\[
a\xi^{3} + 3b\xi^{2}\eta + 3c\xi\eta^{2} + d\eta^{3} = 0, \quad b\xi^{3} + 3c\xi^{2}\eta + 3d\xi\eta^{2} + e\eta^{3} = 0,
\]
identically have a pair of equal roots in common; in other words, that it is possible to satisfy
\[
(A\xi + B\eta)(a\xi^{3} + 3b\xi^{2}\eta + 3c\xi\eta^{2} + d\eta^{3}) + (A'\xi + B'\eta)(b\xi^{3} + 3c\xi^{2}\eta + 3d\xi\eta^{2} + e\eta^{3}) = 0.
\]
We obtain, equating to zero the separate terms of this equation, and eliminating $A$, $B$, $A'$, $B'$,
\[
\begin{vmatrix}
a, & 3b, & 3c, & d \\
b, & 3c, & 3d, & e \\
a, & 3b, & 3c, & d, & \cdot \\
b, & 3c, & 3d, & e, & \cdot
\end{vmatrix} = 0.
\]
It is not at first sight obvious what connection these equations have with the two, $I = 0$, $J = 0$, but by actual expansion they reduce themselves to the following five,
\begin{align*}
3[2(ce - d^{2})I - 3eJ] &= 0, \\
3[(bc - cd)I - 3dJ] &= 0, \\
[-(ae + 2bd - 3c^{2})I + 9cJ] &= 0, \\
3[(ad - bc)I - 3bJ] &= 0, \\
3[2(ac - b^{2})I - 3aJ] &= 0;
\end{align*}
which are satisfied by $I = 0$, $J = 0$. By the theorem above given, the equations are to be considered as forming a system of the tenth order; the system must therefore be considered as composed of the system $I = 0$, $J = 0$, and of a system of the fourth order. The system of the fourth order may be written in the form
\[
2(ac - b^{2}) : ad - bc : ae + 2bd - 3c^{2} : be - cd : 2(ce - d^{2}) = a : b : 3c : d : e.
\]

\newpage

\setcounter{page}{461}

but to justify this, it must be shown first that these equations reduce themselves to two independent equations; and next that system is really one of the fourth order.

We may remark in the first place, that if
\[
a\xi^{4} + 4b\xi^{3}\eta + 6c\xi^{2}\eta^{2} + 4d\xi\eta^{3} + e\eta^{4}
\]
is a perfect square, the coefficients will be proportional to those of $(\alpha\xi^{2} + 2\beta\xi\eta + \gamma\eta^{2})^{2}$. Thus the conditions requisite in order that $u$ may be a perfect square, are given by the system
\[
2(ac - b^{2}) : (ad - bc) : ae + 2bd - 3c^{2} : be - cd : 2(ce - d^{2}) = a : b : 3c : d : e,
\]
or these equations are equivalent to two independent equations only (this may be easily verified \textit{a posteriori}); and by writing $3J$ in the form
\[
e(ac - b^{2}) - 2d(ad - bc) + c(ae + 2bd - 3c^{2}) - 2b(be - cd) + a(ce - d^{2}),
\]
the remaining equations of the complete system (3) are immediately deduced; thus the latter system contains only two independent equations. (The preceding reasoning shows that the system (3) expresses the conditions in order that the equations
\[
a\xi^{3} + 3b\xi^{2}\eta + 3c\xi\eta^{2} + d\eta^{3} = 0, \quad b\xi^{3} + 3c\xi^{2}\eta + 3d\xi\eta^{2} + e\eta^{3} = 0,
\]
may have a pair of unequal roots in common: we have already seen that the equations $I = 0$, $J = 0$ represent the conditions in order that these two equations may have a pair of equal roots in common.) Finally, to verify \textit{a posteriori} the fact of the system (3) being one only of the fourth order, we may, as Mr Salmon has done in the memoir above referred to, represent the system by the two equations
\[
a(ce - d^{2}) - e(ac - b^{2}) = 0, \quad e(ad - bc) - 2b(ce - d^{2}) = 0,
\]
that is, by $ad^{2} - eb^{2} = 0$, $2bd^{2} - 3ce + ae = 0$.

These equations contain the extraneous system $(a = 0, b = 0)$ and the extraneous systems $(b = 0, d = 0)$ and $(d = 0, e = 0)$, each of which last, as Mr Salmon has remarked from geometrical considerations, counts double, or the system is one of the $4^{\text{th}}$ order only.\footnote{More generally whatever be the order of $u$, if $u$ contain a square factor, this square factor may easily be shown to occur in $\frac{d^{2}u}{d\xi^{2}} \cdot \frac{d^{2}u}{d\eta^{2}} - \left(\frac{d^{2}u}{d\xi\, d\eta}\right)^{2}$.}



\newpage
% ============================================================
% PAPER 78
% ============================================================
\section*{78. Note on the motion of rotation of a solid of revolution.}
\markboth{Note on the motion of rotation of a solid of revolution.}{}
\addcontentsline{toc}{section}{78. Note on the motion of rotation of a solid of revolution.}

\begin{flushright}
\textit{[From the Cambridge and Dublin Mathematical Journal, vol.\ iv.\ (1849), pp.\ 268--270.]}
\end{flushright}

\setcounter{page}{462}

\textsc{Using} the notation employed in my former papers on the subject of rotation (Cambridge Math.\ Journal, vol.~III.\ pp.~224--232, [6]; and Cambridge and Dublin Math.\ Journal, vol.~I.\ pp.~167--264, [37]), suppose $B = A$, then $r$ is constant, equal $n$ suppose; and writing
\[
\theta = \nu t + \tau,
\]
also putting
\[
\epsilon = \nu t + \tau,
\]
(where $\tau$ is an arbitrary constant) the values of $p$, $q$, $r$ are easily seen to be given by the equations
\begin{align*}
p &= M\sin\theta, \\
q &= M\cos\theta, \\
r &= n,
\end{align*}
(where $M$ is arbitrary). And consequently
\[
h = A[M^{2} + n(n - \nu)], \qquad k^{2} = A^{2}[M^{2} + (n - \nu)^{2}].
\]
Also, since $a^{2} + b^{2} + c^{2} = k^{2}$, we may write
\begin{align*}
a &= -k\sin i\cos j, \\
b &= k\cos i\cos j, \\
c &= k\sin j,
\end{align*}
$k$ having the value above given, and the angles $i$, $j$ being arbitrary.

\newpage

\setcounter{page}{463}

From the equations (12) and (16) in the second of the papers quoted, we deduce
\begin{align*}
2\nu &= k\{k + A(n - \nu)\sin j + MA\cos j\cos(\theta + i)\}, \\
\lambda &= k\{n\sin j + M\cos j\cos(\theta + i)\}, \\
\mu &= -\nu kMA\cos j\sin(\theta + i),
\end{align*}
(values which verify as they should do the equation (19)). Hence, from the equation (27), writing $\frac{d\psi}{dt} = \nu = -\frac{d\theta}{dt}$, we have
\[
2\tan^{-1}\frac{\mu}{\lambda} = \theta + \frac{1}{k}\int \frac{k + A(n - \nu)\sin j + MA\cos j\cos(\theta + i)}{k + A(n - \nu)\sin j + MA\cos j\cos(\theta + i)}\, d\theta.
\]
This is easily integrated; but the only case which appears likely to give a simple result is when the quantity under the integral sign is constant, or
\[
A(k + kn\sin j) = k[k + A(n - \nu)\sin j],
\]
or
\[
Ah - k^{2} + Ak\nu\sin j = 0;
\]
that is,
\[
A(n - \nu) + k\sin j = 0:
\]
whence
\[
\sin j = \frac{A(\nu - n)}{k}, \quad \cos j = \frac{AM}{k}, \quad \text{or } \tan j = \frac{-(n - \nu)}{M} = \frac{-\frac{(n - \nu)}{n}}{\frac{M}{n}}.
\]
Observing that $\frac{\pi}{2} - j$ is the inclination of the axis of $z$ to the normal to the invariable plane, this equation shows that the supposition above is not any restriction upon the generality of the motion, but amounts only to supposing that the axis of $z$ (which is a line fixed in space) is taken upon the surface of a certain right cone having for its axis the perpendicular to the invariable plane. Resuming the solution of the problem, we have
\[
2\tan^{-1}\frac{\mu}{\lambda} = \theta + \frac{k}{\nu A}\theta,
\]
which may also be written under the form
\[
2\tan^{-1}\frac{\mu}{\lambda} = \theta_{1} + 2\epsilon,
\]
(where $\theta_{1} = \theta + \frac{\tau}{j}$). And hence
\[
\Omega = k\tan\tfrac{1}{2}(\theta_{1} + 2\epsilon) - k\tan\tfrac{1}{2}\theta_{1},
\]
where

\newpage

\setcounter{page}{464}

Substituting these values,
\begin{align*}
a &= -MA\sin i, \quad b = MA\cos i, \quad c = -A(n - \nu), \\
\omega &= A^{2}M^{2}\{1 + \cos(\theta + i)\} = 2M^{2}A^{2}\cos^{2}\tfrac{1}{2}(\theta + i);
\end{align*}
and substituting in the equations (14) the values of $\lambda$, $\mu$, $\nu$ reduce themselves to
\begin{align*}
\lambda &= \frac{1}{MA\cos\tfrac{1}{2}(\theta - i)}[k\tan\tfrac{1}{2}\epsilon\sin\tfrac{1}{2}(\theta - i) - A(n - \nu)\cos\tfrac{1}{2}(\theta - i)], \\
\mu &= \frac{1}{MA\cos\tfrac{1}{2}(\theta - i)}[k\tan\tfrac{1}{2}\epsilon\cos\tfrac{1}{2}(\theta - i) + A(n - \nu)\sin\tfrac{1}{2}(\theta - i)], \\
\nu &= \tan\tfrac{1}{2}(\theta + i);\ kt,
\end{align*}
where, recapitulating, $\theta = \nu t + \tau$, $2\epsilon = \frac{\tau}{j} + \theta_{1}$.

I may notice, in connexion with the problem of rotation, a memoir, ``Specimen Inaugurale de motu gyratorio corporis rigidi \&c.,'' by A.\ S.\ Rueb (Utrecht, 1834), which contains some very interesting developments of the ordinary solution of the problem, by means of the theory of elliptic functions.



\newpage
% ============================================================
% PAPER 79
% ============================================================
\section*{79. On a system of equations connected with Malfatti's problem, and on another algebraical system.}
\markboth{On a system of equations connected with Malfatti's problem}{}
\addcontentsline{toc}{section}{79. On a system of equations connected with Malfatti's problem, and on another algebraical system.}

\begin{flushright}
\textit{[From the Cambridge and Dublin Mathematical Journal, vol.\ iv.\ (1849), pp.\ 270--275.]}
\end{flushright}

\setcounter{page}{465}

\textsc{Consider} the equations
\begin{align*}
by^{2} + cz^{2} + 2fyz &= \theta\phi' a(bc - f^{2}), \\
cz^{2} + ax^{2} + 2gzx &= \phi\psi' b(ca - g^{2}), \\
ax^{2} + by^{2} + 2hxy &= \psi\theta' c(ab - h^{2});
\end{align*}
or, as they may be more conveniently written,
\begin{align*}
by^{2} + cz^{2} + 2fyz &= \theta\phi' a\mathfrak{A}, \\
cz^{2} + ax^{2} + 2gzx &= \phi\psi' b\mathfrak{B}, \\
ax^{2} + by^{2} + 2hxy &= \psi\theta' c\mathfrak{C}.
\end{align*}
The second and third equations give
\[
(g^{2} - h^{2})x^{2} - b^{2}y^{2} + c^{2}z^{2} + 2cg\mathfrak{C}zx - 2bh\mathfrak{B}xy = 0,
\]
hence $[(g^{2} - h^{2})x - bh\mathfrak{B}y + cg\mathfrak{C}z]^{2} - \mathfrak{B}\mathfrak{C}(-bgy + chz)^{2} = 0$, and consequently
\[
(\mathfrak{A}\alpha - \gamma^{2})x - b\sqrt{\mathfrak{B}\mathfrak{C}}(\sqrt{\mathfrak{B}\mathfrak{C}} + \sqrt{\mathfrak{A}\mathfrak{B}})y + c\sqrt{\mathfrak{A}\mathfrak{C}}(\sqrt{\mathfrak{B}\mathfrak{C}} + h\sqrt{\mathfrak{A}\mathfrak{B}})z = 0:
\]
dividing this by $g\sqrt{\mathfrak{B}\mathfrak{C}} + h\sqrt{\mathfrak{A}\mathfrak{B}}$, and writing down the system of equations to which the equation thus obtained belongs,
\begin{align*}
(\sqrt{\mathfrak{B}\mathfrak{C}} - h\sqrt{\mathfrak{A}\mathfrak{B}})x - b\sqrt{\mathfrak{A}\mathfrak{C}}\,y + c\sqrt{\mathfrak{A}\mathfrak{B}}\,z &= 0, \\
a\sqrt{\mathfrak{B}\mathfrak{C}}\,x + (\sqrt{\mathfrak{C}\mathfrak{A}} - f\sqrt{\mathfrak{B}\mathfrak{C}})y - c\sqrt{\mathfrak{A}\mathfrak{B}}\,z &= 0, \\
-a\sqrt{\mathfrak{B}\mathfrak{C}}\,x + b\sqrt{\mathfrak{C}\mathfrak{A}}\,y + (\sqrt{\mathfrak{A}\mathfrak{B}} - g\sqrt{\mathfrak{C}\mathfrak{A}})z &= 0.
\end{align*}

\newpage

\setcounter{page}{466}

Hence also
\begin{align*}
(\sqrt{\mathfrak{A}\mathfrak{B}} + b\sqrt{\mathfrak{B}\mathfrak{C}} - f\sqrt{\mathfrak{C}\mathfrak{A}})x - (\sqrt{\mathfrak{B}\mathfrak{C}} - f\sqrt{\mathfrak{C}\mathfrak{A}} + c\sqrt{\mathfrak{A}\mathfrak{B}})z &= 0, \\
(\sqrt{\mathfrak{A}\mathfrak{B}} - h\sqrt{\mathfrak{B}\mathfrak{C}} + b\sqrt{\mathfrak{C}\mathfrak{A}})x - (\sqrt{\mathfrak{A}\mathfrak{B}} + b\sqrt{\mathfrak{B}\mathfrak{C}} + f\sqrt{\mathfrak{C}\mathfrak{A}})y &= 0;
\end{align*}
these equations may be written
\begin{align*}
\tfrac{1}{2}\{\mathfrak{F}\sqrt{\mathfrak{A}} - \mathfrak{G}\sqrt{\mathfrak{B}} - \mathfrak{H}\sqrt{\mathfrak{C}} + \sqrt{(\mathfrak{A}\mathfrak{B}\mathfrak{C})}\}\{[\mathfrak{F} + \sqrt{(\mathfrak{B}\mathfrak{C})}]y - [\mathfrak{G} + \sqrt{(\mathfrak{C}\mathfrak{A})}]z\} &= 0, \\
\tfrac{1}{2}\{-\mathfrak{F}\sqrt{\mathfrak{A}} + \mathfrak{G}\sqrt{\mathfrak{B}} - \mathfrak{H}\sqrt{\mathfrak{C}} + \sqrt{(\mathfrak{A}\mathfrak{B}\mathfrak{C})}\}\{[\mathfrak{G} + \sqrt{(\mathfrak{C}\mathfrak{A})}]z - [\mathfrak{H} + \sqrt{(\mathfrak{A}\mathfrak{B})}]x\} &= 0, \\
\tfrac{1}{2}\{-\mathfrak{F}\sqrt{\mathfrak{A}} - \mathfrak{G}\sqrt{\mathfrak{B}} + \mathfrak{H}\sqrt{\mathfrak{C}} + \sqrt{(\mathfrak{A}\mathfrak{B}\mathfrak{C})}\}\{[\mathfrak{F} + \sqrt{(\mathfrak{B}\mathfrak{C})}]x - [\mathfrak{H} + \sqrt{(\mathfrak{A}\mathfrak{B})}]y\} &= 0;
\end{align*}
where, as usual,
\[
\mathfrak{F} = gh - af, \quad \mathfrak{G} = hf - bg, \quad \mathfrak{H} = fg - ch, \quad K = abc - af^{2} - bg^{2} - ch^{2} + 2fgh;
\]
(in fact the coefficient of $y$ in the first equation is
\[
\tfrac{1}{2}\{(\mathfrak{F}\mathfrak{G} - \mathfrak{H}^{2})\sqrt{\mathfrak{A}} + (ac - \mathfrak{G}^{2})\sqrt{\mathfrak{B}} - (\mathfrak{C}\mathfrak{F} - a\mathfrak{H})\sqrt{\mathfrak{C}}\} = \mathfrak{F}\sqrt{\mathfrak{A}} + \mathfrak{G}\sqrt{\mathfrak{B}} - f\sqrt{\mathfrak{C}},
\]
as it should be, and similarly for the coefficients of the remaining terms). We have therefore
\[
[\mathfrak{F} + \sqrt{(\mathfrak{B}\mathfrak{C})}]x = [\mathfrak{G} + \sqrt{(\mathfrak{C}\mathfrak{A})}]y = [\mathfrak{H} + \sqrt{(\mathfrak{A}\mathfrak{B})}]z;
\]
or, what comes to the same thing,
\begin{align*}
yz &= \tfrac{1}{2}a\{\mathfrak{F} + \sqrt{(\mathfrak{B}\mathfrak{C})}\}, \\
zx &= \tfrac{1}{2}b\{\mathfrak{G} + \sqrt{(\mathfrak{C}\mathfrak{A})}\}, \\
xy &= \tfrac{1}{2}c\{\mathfrak{H} + \sqrt{(\mathfrak{A}\mathfrak{B})}\}.
\end{align*}
Now
\begin{align*}
a\{\mathfrak{G} + \sqrt{(\mathfrak{C}\mathfrak{A})}\}\{\mathfrak{H} + \sqrt{(\mathfrak{A}\mathfrak{B})}\} &= \{\mathfrak{F} + \sqrt{(\mathfrak{B}\mathfrak{C})}\}\{abc - fgh + f\sqrt{(\mathfrak{B}\mathfrak{C})} - g\sqrt{(\mathfrak{C}\mathfrak{A})} - h\sqrt{(\mathfrak{A}\mathfrak{B})}\}, \\
b\{\mathfrak{H} + \sqrt{(\mathfrak{A}\mathfrak{B})}\}\{\mathfrak{F} + \sqrt{(\mathfrak{B}\mathfrak{C})}\} &= \{\mathfrak{G} + \sqrt{(\mathfrak{C}\mathfrak{A})}\}\{abc - fgh - f\sqrt{(\mathfrak{B}\mathfrak{C})} + g\sqrt{(\mathfrak{C}\mathfrak{A})} - h\sqrt{(\mathfrak{A}\mathfrak{B})}\}, \\
c\{\mathfrak{F} + \sqrt{(\mathfrak{B}\mathfrak{C})}\}\{\mathfrak{G} + \sqrt{(\mathfrak{C}\mathfrak{A})}\} &= \{\mathfrak{H} + \sqrt{(\mathfrak{A}\mathfrak{B})}\}\{abc - fgh - f\sqrt{(\mathfrak{B}\mathfrak{C})} - g\sqrt{(\mathfrak{C}\mathfrak{A})} + h\sqrt{(\mathfrak{A}\mathfrak{B})}\},
\end{align*}
{as readily appears by writing the first of these equations under the form
\[
a\{\mathfrak{G} + \sqrt{(\mathfrak{C}\mathfrak{A})}\}\{\mathfrak{H} + \sqrt{(\mathfrak{A}\mathfrak{B})}\} = \{\mathfrak{F} + \sqrt{(\mathfrak{B}\mathfrak{C})}\}\{aa - f^{2} + f\sqrt{(\mathfrak{B}\mathfrak{C})} - g\sqrt{(\mathfrak{C}\mathfrak{A})} - h\sqrt{(\mathfrak{A}\mathfrak{B})}\},
\]
and comparing the rational term and the coefficients of $\sqrt{(\mathfrak{B}\mathfrak{C})}$, $\sqrt{(\mathfrak{C}\mathfrak{A})}$, $\sqrt{(\mathfrak{A}\mathfrak{B})}$}.

\newpage

\setcounter{page}{467}

Hence, observing the values of $yz$, $zx$, $xy$, we find
\begin{align*}
x^{2} &= \tfrac{1}{2a}\{abc - fgh + f\sqrt{(\mathfrak{B}\mathfrak{C})} - g\sqrt{(\mathfrak{C}\mathfrak{A})} - h\sqrt{(\mathfrak{A}\mathfrak{B})}\}, \\
y^{2} &= \tfrac{1}{2b}\{abc - fgh - f\sqrt{(\mathfrak{B}\mathfrak{C})} + g\sqrt{(\mathfrak{C}\mathfrak{A})} - h\sqrt{(\mathfrak{A}\mathfrak{B})}\}, \\
z^{2} &= \tfrac{1}{2c}\{abc - fgh - f\sqrt{(\mathfrak{B}\mathfrak{C})} - g\sqrt{(\mathfrak{C}\mathfrak{A})} + h\sqrt{(\mathfrak{A}\mathfrak{B})}\}.
\end{align*}
Hence, forming the value of any one of the functions $by^{2} + cz^{2} + 2fyz$, $cz^{2} + ax^{2} + 2gzx$, $ax^{2} + by^{2} + 2hxy$, we obtain $\mathfrak{s} = \theta\phi\psi$; or we have
\begin{align*}
x^{2} &= \tfrac{1}{2a}\{abc - fgh + f\sqrt{(\mathfrak{B}\mathfrak{C})} - g\sqrt{(\mathfrak{C}\mathfrak{A})} - h\sqrt{(\mathfrak{A}\mathfrak{B})}\}, \\
y^{2} &= \tfrac{1}{2b}\{abc - fgh - f\sqrt{(\mathfrak{B}\mathfrak{C})} + g\sqrt{(\mathfrak{C}\mathfrak{A})} - h\sqrt{(\mathfrak{A}\mathfrak{B})}\}, \\
z^{2} &= \tfrac{1}{2c}\{abc - fgh - f\sqrt{(\mathfrak{B}\mathfrak{C})} - g\sqrt{(\mathfrak{C}\mathfrak{A})} + h\sqrt{(\mathfrak{A}\mathfrak{B})}\}, \\
2yz &= \tfrac{1}{a}\{\mathfrak{F} + \sqrt{(\mathfrak{B}\mathfrak{C})}\}, \\
2zx &= \tfrac{1}{b}\{\mathfrak{G} + \sqrt{(\mathfrak{C}\mathfrak{A})}\}, \\
2xy &= \tfrac{1}{c}\{\mathfrak{H} + \sqrt{(\mathfrak{A}\mathfrak{B})}\}.
\end{align*}

It may be remarked that the equations
\begin{align*}
hy^{2} + cz^{2} + 2f'yz &= L, \\
c'z^{2} + a'x^{2} + 2g'zx &= M, \\
a''x^{2} + b''y^{2} + 2h''xy &= N,
\end{align*}
in which the coefficients are supposed to be such that the functions
\begin{align*}
M(a''x^{2} + b''y^{2} + 2h''xy) - N(c'z^{2} + a'x^{2} + 2g'yz), \\
N(by^{2} + cz^{2} + 2fyz) - L(a''x^{2} + b''y^{2} + 2h''yz), \\
L(c'z^{2} + a'x^{2} + 2g'yz) - M(by^{2} + cz^{2} + 2fyz),
\end{align*}
are each of them decomposable into linear factors, may always be reduced to a system of equations similar to those which have just been solved.

Suppose
\[
f = g = h = \theta = \phi = \psi = \sqrt{(bc)} = \sqrt{(ca)} = \sqrt{(ab)} = \ldots,
\]
and write $\sqrt{X}$, $\sqrt{Y}$, $\sqrt{Z}$ instead of $x$, $y$, $z$. The equations to be solved become
\begin{align*}
bY + cZ + 2f\sqrt{(YZ)} &= (b + c)r, \\
cZ + aX + 2g\sqrt{(ZX)} &= (c + a)r, \\
aX + bY + 2h\sqrt{(XY)} &= (a + b)r,
\end{align*}

\newpage

\setcounter{page}{468}

where $r^{2} = \frac{abc}{a + b + c}$, and the solution is
\begin{align*}
X &= \tfrac{1}{2a}\{a + b + c - r + \sqrt{(r^{2} + a^{2})} - \sqrt{(r^{2} + b^{2})} - \sqrt{(r^{2} + c^{2})}\}, \\
Y &= \tfrac{1}{2b}\{a + b + c - r - \sqrt{(r^{2} + a^{2})} + \sqrt{(r^{2} + b^{2})} - \sqrt{(r^{2} + c^{2})}\}, \\
Z &= \tfrac{1}{2c}\{a + b + c - r - \sqrt{(r^{2} + a^{2})} - \sqrt{(r^{2} + b^{2})} + \sqrt{(r^{2} + c^{2})}\}, \\
\sqrt{(yz)} &= \tfrac{1}{2}\{r - a + \sqrt{(r^{2} + a^{2})}\}, \\
\sqrt{(xy)} &= \tfrac{1}{2}\{r - c + \sqrt{(r^{2} + c^{2})}\},
\end{align*}
a system of formul\ae\ which contain the solution of the problem ``In a given triangle to inscribe three circles such that each circle touches the remaining two circles and also two sides of the triangle.'' In fact, if $r$ denote the radius of the inscribed circle, and $a$, $b$, $c$ the distances of the angles of the triangle from the points where the sides are touched by the inscribed circle (quantities which it is well known satisfy the condition $r^{2} = \frac{abc}{a+b+c}$), also if $x$, $y$, $z$ denote the radii of the required circles, there is no difficulty whatever in obtaining for the determination of $x$, $y$, $z$, the above system of equations. The problem in question was first proposed and solved by an Italian geometer named Malfatti, and has been called after him Malfatti's problem. His solution, dated 1803, and published in the 10th volume of the Transactions of the Italian Academy of Sciences, appears to have consisted in showing that the values first found for the radii of the three circles satisfy the equations given above, without any indication of the process of obtaining the expressions for these radii. Further information as to the history of the problem may be found in the memoir ``Das Malfattische Problem neu gel\"{o}st von C.\ Adams,'' Winterthur, 1846.

In connexion with the preceding investigations may be considered the problem of determining $l$ and $m$ from the equations
\begin{align*}
B(l + \epsilon)^{2} - 2H(l + \epsilon)m + (A + \lambda)m^{2} &= 0, \\
A'(m + \epsilon)^{2} - 2H(m + \epsilon)l + (B + \lambda)l^{2} &= 0;
\end{align*}
which express that the function
\[
\theta\phi' U + (lx + my)^{2}, \quad (U = Ax^{2} + 2Hxy + By^{2}),
\]
has for one of its factors a factor of $U + x^{2}$, and for the other of its factors a factor of $U + y^{2}$: There is no difficulty in solving these equations; and if we write
\[
K = AB - H^{2}, \quad n\lambda = \sqrt{(-K - B)}, \quad n\mu = \sqrt{(-K - A)},
\]

\end{document}
